\documentclass[12pt,a4paper]{report}
\renewcommand\normalsize{\fontsize{15}{18}\selectfont}
\normalsize
\usepackage{float}
\usepackage{amsmath}
\usepackage{amssymb}
\usepackage{booktabs}
\usepackage{tikz}
\usetikzlibrary{positioning}
\usepackage{pgfplots}
\pgfplotsset{compat=1.18}
\usepackage{tikz}
\usepackage{amsmath}
\usepackage{amssymb}
\usepackage{pgfplots}
\pgfplotsset{compat=1.18}
\usepackage[utf8]{inputenc}
\usepackage[T1]{fontenc}
\usepackage[french]{babel}
\usepackage{graphicx}
\usepackage{changepage}
\usepackage{geometry}
\usepackage{svg}
\usepackage{caption}
\geometry{
  left=1.5cm,
  right=1.5cm,
  top=2cm,
  bottom=2.5cm
}
\usepackage{xcolor}
\usepackage{listings}
% Définition des couleurs
\definecolor{background}{HTML}{1E1E1E} % fond VS Code
\definecolor{keyword}{HTML}{D19A66}    % orange (keywords)
\definecolor{class}{HTML}{56B6C2}      % cyan (class, def)
\definecolor{function}{HTML}{61AFEF}   % bleu clair (fonctions)
\definecolor{string}{HTML}{C678DD}     % violet (strings)
\definecolor{comment}{HTML}{7F848E}    % gris (commentaires)
\definecolor{type}{HTML}{98C379}       % vert clair (types)
\definecolor{defaulttext}{HTML}{ECECEC} % texte par défaut

\lstset{
  language=Python,
  backgroundcolor=\color{background},
  basicstyle=\ttfamily\footnotesize\color{defaulttext},
  keywordstyle=\color{keyword}\bfseries,
  stringstyle=\color{string},
  commentstyle=\color{comment}\itshape,
  identifierstyle=\color{defaulttext},
  emph={self, np, List}, emphstyle=\color{type},
  emph={[2]class,def}, emphstyle={[2]\color{class}\bfseries},
  emph={[3]return}, emphstyle={[3]\color{function}\bfseries},
  showstringspaces=false,
  breaklines=true,
  tabsize=4,
  % >>> suppression des espaces verticaux entre lignes
  aboveskip=0pt,
  belowskip=0pt,
  lineskip=-1pt
}

\usepackage{xcolor}
\usepackage{lipsum} 
\usepackage{titlesec}
\usepackage{fancyhdr}
\usepackage{setspace}
\usepackage{titling}
\usepackage{multicol}
\usepackage{array}
\usepackage{lmodern}
\usepackage{parskip}
\usepackage{caption}
\usepackage{tocloft} % Pour améliorer la table des matières
\usepackage{titlesec}
\titleclass{\subsubsubsection}{straight}[\subsubsection]

\usepackage{etoolbox}
\usepackage{subcaption}

% Pour afficher ce niveau dans la table des matières
\setcounter{tocdepth}{4}

% Style global fancyhdr
\pagestyle{fancy}
\fancyhf{}  % reset tout

% En-tête : logos à gauche et droite, taille adaptée
\fancyhead[L]{%
  \includegraphics[height=1.5cm,keepaspectratio]{Images_/logo_ensam_mek.png}
}
\fancyhead[R]{%
  % Pense à ajouter le logo de AI Movement dans ton dossier Images_
  \includegraphics[height=2cm,keepaspectratio]{Images_/aimovement_logo.png}
}
\renewcommand{\headrulewidth}{0.4pt} % ligne sous en-tête

% Pied de page : nom à gauche, année au centre, numéro à droite
\fancyfoot[L]{Amine Essahraoui}
\fancyfoot[C]{2025 / 2026}
\fancyfoot[R]{\thepage}
\renewcommand{\footrulewidth}{0.4pt} % ligne au-dessus pied

% Ha ch7al f toul dyal l'en-tête (bach thaz les logos)
\setlength{\headheight}{1.5cm} 

% Ha l'espace bin l'en-tête (lkhat) w lktba
\setlength{\headsep}{0.9cm}

% Début du document
\begin{document}

\begin{titlepage}
  \thispagestyle{empty}

  % En-tête avec les logos (ENSAM à gauche, AI Movement à droite)
  \begin{minipage}[t]{0.48\textwidth}
    \begin{flushleft}
      \includegraphics[height=1.5cm]{Images_/logo_ensam_mek.png}\\[0.2cm]
      {\small \textbf{École Nationale Supérieure d'Arts}\\
      \textbf{et Métiers - Meknès}}
    \end{flushleft}
  \end{minipage}
  \hfill
  \begin{minipage}[t]{0.48\textwidth}
    \begin{flushright}
      \includegraphics[height=2.3cm]{Images_/aimovement_logo.png}\\[0.2cm]
      {\small \textbf{AI Movement - UM6P}\\
      Salé Al Jadida}
    \end{flushright}
  \end{minipage}

  \vspace*{3.5cm}

  \begin{center}
    % Type de document
    {\LARGE \textbf{Rapport de Projet de Fin D'année}}\\[1.5cm]
    
    % Titre du projet avec des lignes esthétiques
    \rule{\linewidth}{0.5pt}\\[0.6cm]
    {\Huge \linespread{1.1}\selectfont \textbf{Conception, Optimisation et Évaluation des performances d'un Pipeline à Haute Performance pour l'Entraînement Distribué de Modèles Vision-Langage}\par}
    \vspace{0.4cm}
    \rule{\linewidth}{0.5pt}\\[1.5cm]

    % Filière
    {\large Filière : Intelligence Artificielle et Technologies des Données}\\[2.5cm]

    \begin{large}
      \textbf{Réalisé par :} Amine Essahraoui \\[0.5cm]
      \textbf{Encadrant professionnel :} M. Youssef Hmamouche \\[0.5cm]
      % \textbf{Encadrant académique :}
    \end{large}

    \vfill

    % Année universitaire
    {\large Année universitaire : 2025 -- 2026}
  \end{center}
\end{titlepage}

\newpage
\thispagestyle{empty}
~
\newpage

\begin{adjustwidth}{-1cm}{-1cm}

\begin{center}
    {\LARGE \textbf{Remerciements}}
\end{center}

\vspace{1.5cm}

\begin{center}
\begin{minipage}{0.9\textwidth}
\large

Je tiens à exprimer mes sincères remerciements à l'École Nationale Supérieure d'Arts et Métiers de Meknès (ENSAM Meknès) pour la qualité de la formation dispensée.

\vspace{0.7cm}

Je remercie particulièrement les professeurs de la filière Intelligence Artificielle et Technologies des Données (IATD), \textbf{Tawfik Masrour}, \textbf{Hosni Mohammed} et \textbf{Tarik Hajji}, pour leur accompagnement et leurs précieux conseils tout au long de mon parcours académique.

\vspace{0.7cm}

Je souhaite également adresser un grand remerciement à mon encadrant professionnel, Monsieur \textbf{Youssef Hmamouche}, ainsi qu'à toute l'équipe du Centre International d'Intelligence Artificielle du Maroc (\textbf{AI Movement - UM6P}) pour leur accueil chaleureux, leur soutien scientifique constant et l'opportunité qu'ils m'ont offerte d'effectuer ce projet de recherche enrichissant.

\vspace{0.7cm}

Cette expérience m'a permis de mettre en pratique mes connaissances en ingénierie logicielle et en architecture des systèmes distribués, de développer mes compétences dans le domaine du calcul haute performance (HPC) et de contribuer activement à l'optimisation des fondations de l'intelligence artificielle.

\vspace{0.7cm}

Enfin, je remercie toutes les personnes qui ont, de près ou de loin, contribué à la réussite de ce stage et à la rédaction de ce rapport.

\end{minipage}
\end{center}

\vspace{1.5cm}

\begin{flushright}
\textit{Amine Essahraoui}
\end{flushright}

\end{adjustwidth}
\newpage
\thispagestyle{empty}

\begin{center}
    \LARGE \textbf{Liste des abréviations}
\end{center}
\addcontentsline{toc}{section}{Liste des abréviations}
\vspace{1cm}

\begin{tabular}{>{\bfseries}ll}
AVX    & Advanced Vector Extensions \\
CPU    & Central Processing Unit (Processeur Central) \\
CRC    & Cyclic Redundancy Check (Contrôle de Redondance Cyclique) \\
DDP    & Distributed Data Parallel (Parallélisme de Données Distribué) \\
ENSAM  & École Nationale Supérieure d’Arts et Métiers \\
FLOPs  & Floating Point Operations Per Second (Opérations en virgule flottante par seconde) \\
GIL    & Global Interpreter Lock (Verrou Global de l'Interpréteur Python) \\
GPU    & Graphics Processing Unit (Processeur Graphique) \\
H2D    & Host-to-Device (Transfert de l'hôte vers le périphérique) \\
HPC    & High-Performance Computing (Calcul Haute Performance) \\
IA     & Intelligence Artificielle \\
IATD   & Intelligence Artificielle et Technologies des Données \\
I/O    & Input/Output (Entrées/Sorties) \\
ISO    & International Organization for Standardization (Organisation Internationale de Normalisation) \\
JIT    & Just-In-Time (Compilation à la volée) \\
JSON   & JavaScript Object Notation \\
LLM    & Large Language Model (Grand Modèle de Langage) \\
LSB    & Least Significant Bit (Bit de poids faible) \\
LZ4    & Lempel-Ziv 4 (Algorithme de compression sans perte) \\
MESI   & Modified, Exclusive, Shared, Invalid (Protocole de cohérence de cache) \\
mmap   & Memory-Mapped File (Fichier Mappé en Mémoire) \\
PCIe   & Peripheral Component Interconnect Express \\
POSIX  & Portable Operating System Interface \\
RAII   & Resource Acquisition Is Initialization (Acquisition d'une ressource à l'initialisation) \\
RAM    & Random Access Memory (Mémoire à Accès Aléatoire) \\
RNG    & Random Number Generator (Générateur de Nombres Aléatoires) \\
RQ     & Research Question (Question de Recherche) \\
RVB    & Rouge, Vert, Bleu (Espace Colorimétrique) \\
SIMD   & Single Instruction Multiple Data \\
SPSC   & Single Producer Single Consumer (Producteur Unique, Consommateur Unique) \\
TLB    & Translation Lookaside Buffer (Tampon de traduction) \\
UM6P   & Université Mohammed VI Polytechnique \\
ViT    & Vision Transformer \\
VLM    & Vision-Language Model (Modèle Vision-Langage) \\
VQA    & Visual Question Answering (Question-Réponse Visuelle) \\
YAML   & YAML Ain't Markup Language \\
\end{tabular}


\addcontentsline{toc}{section}{Table des figures}
\newpage
\thispagestyle{empty}
\listoffigures

\addcontentsline{toc}{section}{Liste des tableaux}
\newpage
\thispagestyle{empty}
\listoftables


\newpage

% ---------- Début du contenu principal ----------
\pagenumbering{arabic}
\setcounter{page}{1}

\section*{Introduction Générale}
\addcontentsline{toc}{section}{Introduction générale}

\subsection{Contexte et Problématique}

Le présent projet de fin d'études s'inscrit dans le cadre de mes travaux de recherche au sein du Centre International d'Intelligence Artificielle du Maroc (\textbf{AI Movement}), hébergé par l'Université Mohammed VI Polytechnique (\textbf{UM6P}). L'essor fulgurant de l'intelligence artificielle générative repose aujourd'hui sur l'entraînement de modèles fondationnels, notamment les modèles Vision-Langage (VLM), qui nécessitent l'ingestion de pétaoctets de données multimodales (images et textes).

Cependant, cette évolution fait face à une barrière architecturale majeure : le \textit{Data Loading Bottleneck} (goulot d'étranglement lié au chargement des données). Les accélérateurs matériels modernes (GPU) calculent beaucoup plus vite que les processeurs (CPU) ne peuvent lire, décoder et prétraiter les données depuis le disque. Ce déséquilibre engendre un phénomène critique appelé \textit{GPU Starvation} (famine du GPU), où le processeur graphique reste inactif en attendant que le prochain lot de données (\textit{batch}) soit prêt, limitant ainsi drastiquement l'efficacité et la scalabilité des entraînements distribués.

\subsection{Objectifs du Projet}

Face aux limites des solutions standards telles que le \texttt{DataLoader} natif de PyTorch, l'objectif principal de ce projet est la conception et l'implémentation d'un \textbf{pipeline de chargement de données multimodales distribué et à très haute performance (HPC)}. 

Pour atteindre cet objectif, plusieurs défis techniques ont été relevés :
\begin{itemize}
  \item \textbf{Optimisation des Entrées/Sorties (I/O) :} Élimination des copies mémoire inutiles (Zero-Copy) grâce à l'utilisation d'un chargeur asynchrone en C++ basé sur la projection mémoire (\texttt{mmap}).
  \item \textbf{Réduction de la complexité algorithmique :} Implémentation d'une stratégie de \textit{Dynamic Padding} pour éviter le gaspillage de calculs (FLOPs) induit par le \textit{Static Padding} dans l'attention quadratique des Transformers.
  \item \textbf{Scalabilité Multiprocessing :} Développement d'un pipeline Python en flux continu (\textit{streaming}) à mémoire bornée pour le prétraitement des données.
  \item \textbf{Entraînement Distribué :} Création d'un système de partitionnement déterministe (\textit{ShardManifest}) garantissant l'intégrité des données dans un environnement multi-nœuds (DDP).
\end{itemize}

\subsection{Organisation du Rapport}

Ce rapport retrace la démarche scientifique et technique adoptée pour construire cette architecture. Il s’articule autour des chapitres suivants :

\begin{itemize}
  \item Le \textbf{Chapitre 1} présente l'organisme d'accueil (AI Movement - UM6P) et pose le cadre général du projet.
  \item Le \textbf{Chapitre 2} expose le cadre théorique et les fondements mathématiques (modèle Roofline, théorie de mmap, complexité temporelle de l'attention) justifiant nos choix architecturaux.
  \item Le \textbf{Chapitre 3} détaille la conception et la modélisation de l'architecture globale en trois étages.
  \item Le \textbf{Chapitre 4} est consacré à l'implémentation du prétraitement Python, de la sérialisation binaire et du manifeste distribué.
  \item Le \textbf{Chapitre 5} plonge dans les aspects bas niveau avec l'implémentation du runtime C++, intégrant les structures de données sans verrou (lock-free) et le transfert GPU.
  \item Enfin, le \textbf{Chapitre 6} valide scientifiquement notre solution à travers une suite exhaustive de benchmarks comparatifs.
\end{itemize}

% PAGE DE CHAPITRE STYLISÉE AVEC LOGOS ET TITRE CENTRÉ
\begin{titlepage}
  \thispagestyle{empty} % Supprime les en-têtes/pieds pour cette page

  % LOGOS HAUT DE PAGE
  \begin{minipage}{0.48\textwidth}
  \includegraphics[width=7cm]{Images_/logo_ensam_mek.png}\\[0.2cm]
  \textbf{}
\end{minipage}
\hfill
\begin{minipage}{0.48\textwidth}
  \raggedleft
  % Pense à ajouter le logo de AI Movement dans ton dossier Images_
  \includegraphics[width=5.5cm]{Images_/aimovement_logo.png}\\[0.4cm]
  \textbf{} \\
\end{minipage}

  % ESPACE VIDE POUR PLACEMENT VERTICAL AU CENTRE
  \vspace*{5cm}

  % TITRE AU CENTRE
  \begin{center}
    \rule{\linewidth}{0.5pt}\vspace{0.7cm}

    {\Huge\bfseries\color{blue}
    Chapitre 1 :\\[0.5em]
    Présentation de l'organisme d'accueil et contexte du projet}

    \vspace{0.7cm}
    \rule{\linewidth}{0.5pt}
  \end{center}

  \vfill

\end{titlepage}

\chapter{Présentation de l'organisme d'accueil et contexte du projet}

\section{Présentation de l'organisme d'accueil}

\subsection{L'Université Mohammed VI Polytechnique (UM6P)}

L'Université Mohammed VI Polytechnique (UM6P) est une université marocaine à vocation internationale, orientée vers la recherche scientifique, l'innovation, la formation d'excellence et le développement de solutions répondant aux enjeux économiques, technologiques, environnementaux et sociaux du Maroc et du continent africain.

Fondée sous l'impulsion du Groupe OCP, l'UM6P s'inscrit dans une démarche visant à rapprocher la formation académique, la recherche scientifique, le monde industriel et l'entrepreneuriat. Son positionnement repose notamment sur une approche expérimentale de l'apprentissage et de la recherche, ainsi que sur le développement de partenariats avec des institutions académiques, industrielles et publiques aux niveaux national et international.

L'université développe ainsi un modèle fortement orienté vers la production de connaissances et leur transformation en applications concrètes. Cette orientation se traduit par l'importance accordée aux sciences et technologies, à l'intelligence artificielle, à la science des données, à l'agriculture, à l'énergie, à la santé, aux sciences des matériaux, aux sciences sociales et aux disciplines liées au développement durable.

L'UM6P dispose de plusieurs implantations au Maroc et à l'international. Son campus principal se situe à Benguérir, au sein de la Ville Verte Mohammed VI, tandis que d'autres infrastructures universitaires et scientifiques sont notamment implantées à Rabat. Cette organisation permet à l'université de développer un réseau académique et scientifique dépassant le cadre national et de renforcer ses collaborations avec des établissements de recherche et des partenaires industriels internationaux.

\begin{figure}[htbp]
    \centering
    \includegraphics[width=0.65\textwidth]{Images_/logo_um6p.png}
    \caption{Logo de l'Université Mohammed VI Polytechnique (UM6P)}
    \label{fig:logo_um6p}
\end{figure}

L'orientation de l'UM6P vers la recherche appliquée constitue un élément central de son identité. L'université met en place des infrastructures scientifiques et technologiques permettant aux chercheurs, enseignants, doctorants et étudiants de travailler sur des problématiques complexes et multidisciplinaires. Son modèle favorise notamment les interactions entre la recherche fondamentale, l'expérimentation, le développement technologique et les besoins du monde socio-économique.

Cette philosophie se manifeste également à travers le développement de plateformes expérimentales et de laboratoires spécialisés. L'université dispose notamment d'infrastructures consacrées à l'agriculture expérimentale, à l'énergie, aux sciences des matériaux, aux technologies numériques et au calcul scientifique.

Dans ce contexte, l'intelligence artificielle occupe une place stratégique dans l'écosystème scientifique de l'université. Elle constitue à la fois un domaine de recherche transversal et un levier technologique susceptible d'être appliqué à des problématiques liées à l'industrie, à l'agriculture, à l'énergie, à la santé, à l'environnement, à la finance et à de nombreux autres secteurs. C'est dans cette dynamique que s'inscrit le Centre International d'Intelligence Artificielle du Maroc, dénommé \textit{AI Movement}.

\subsection{Le Centre AI Movement}

\subsubsection{Présentation générale}

Le Centre International d'Intelligence Artificielle du Maroc, connu sous le nom d'\textit{AI Movement}, constitue un centre d'excellence consacré à l'intelligence artificielle et aux sciences des données au sein de l'écosystème de l'Université Mohammed VI Polytechnique. Sa création répond à une volonté de développer une capacité nationale et africaine en matière d'intelligence artificielle, tout en favorisant l'ouverture du Maroc sur les réseaux scientifiques et technologiques internationaux.

AI Movement se positionne comme un pôle dédié à la recherche, à l'innovation et à la transformation par l'intelligence artificielle. Sa vocation ne se limite pas au développement de modèles ou d'outils d'intelligence artificielle : le centre cherche également à favoriser l'émergence d'un écosystème réunissant chercheurs, étudiants, industriels, institutions publiques, entrepreneurs et partenaires internationaux.

Le centre vise notamment à favoriser l'émergence d'un savoir-faire marocain en intelligence artificielle et en sciences des données, à contribuer au positionnement du Maroc comme un pôle régional de l'IA et à accompagner les transformations scientifiques, technologiques, économiques et sociétales induites par le développement de ces technologies.

Le positionnement d'AI Movement s'inscrit ainsi dans une perspective à la fois nationale, africaine et internationale. Cette dimension revêt une importance particulière dans le contexte actuel de développement rapide de l'intelligence artificielle, caractérisé par l'émergence de nouveaux paradigmes d'apprentissage automatique, de modèles génératifs, de systèmes autonomes, d'agents intelligents et de systèmes capables de traiter simultanément plusieurs modalités de données.

\subsubsection{Origine et développement du centre}

Le projet AI Movement a notamment été porté par la Professeure Amal El Fallah Seghrouchni, chercheuse spécialisée en intelligence artificielle distribuée et en systèmes multi-agents. Ses travaux scientifiques portent notamment sur les agents intelligents, les systèmes multi-agents, leur coordination ainsi que différentes applications de l'intelligence artificielle.

La création et le développement du centre s'inscrivent dans une volonté de construire un pont scientifique entre le Maroc, l'Europe et, plus largement, le continent africain. Dans cette perspective, AI Movement a été intégré à l'écosystème de l'UM6P afin de bénéficier de l'environnement académique, scientifique et technologique de l'université tout en développant des collaborations internationales.

La Professeure Amal El Fallah Seghrouchni assure la direction du centre. Sous son impulsion, AI Movement développe des activités combinant recherche scientifique, formation, innovation technologique, transfert industriel et réflexion sur les impacts de l'intelligence artificielle. Le centre constitue ainsi un espace dans lequel les problématiques scientifiques peuvent être étudiées conjointement avec leurs applications et leurs implications sociétales.

\begin{figure}[htbp]
    \centering
    \includegraphics[width=0.55\textwidth]{Images_/logo_ai_movement.png}
    \caption{Logo du Centre International d'Intelligence Artificielle du Maroc --- AI Movement}
    \label{fig:logo_ai_movement}
\end{figure}

\subsubsection{Le Dôme de l'IA}

L'une des infrastructures les plus emblématiques associées à AI Movement est le bâtiment surnommé le \textit{Dôme de l'IA}. Cette infrastructure constitue un espace consacré aux activités scientifiques, académiques et institutionnelles liées à l'intelligence artificielle.

Le Dôme participe à l'identité visuelle et institutionnelle du centre et symbolise la volonté de l'UM6P de disposer d'une infrastructure spécifiquement dédiée à la recherche et à l'innovation dans le domaine de l'intelligence artificielle.

\begin{figure}[htbp]
    \centering
    \includegraphics[width=0.7\textwidth]{Images_/dome_ai_movement.jpg}
    \caption{Dôme du Centre International d'Intelligence Artificielle du Maroc --- AI Movement}
    \label{fig:dome_ai_movement}
\end{figure}

Comme illustré dans la figure~\ref{fig:dome_ai_movement}, l'architecture distinctive du bâtiment contribue à son identification comme un lieu dédié aux activités de recherche, d'innovation et d'échange autour de l'intelligence artificielle. Le Dôme accueille notamment des rencontres scientifiques, des séminaires, des présentations de travaux de recherche et différentes manifestations réunissant chercheurs, étudiants, industriels et partenaires institutionnels.

Les activités scientifiques organisées au sein du centre couvrent notamment des problématiques liées aux systèmes multi-agents, à l'apprentissage automatique, au traitement automatique des langues, à la traduction automatique, aux systèmes autonomes et à diverses applications de l'intelligence artificielle.

\subsubsection{L'environnement de calcul et le supercalculateur Toubkal}

Le développement de travaux de recherche en intelligence artificielle nécessite des capacités de calcul importantes, notamment pour l'entraînement de modèles d'apprentissage profond, le traitement de grands volumes de données, l'expérimentation de modèles complexes et l'exécution de simulations numériques. L'écosystème de l'UM6P dispose à ce titre d'une infrastructure majeure de calcul haute performance : le supercalculateur \textit{Toubkal}.

Il convient de distinguer le centre AI Movement de l'infrastructure Toubkal : Toubkal constitue une infrastructure de calcul haute performance de l'UM6P, hébergée à Benguérir. Elle peut néanmoins soutenir des travaux scientifiques nécessitant des ressources de calcul intensif, notamment dans le domaine de l'intelligence artificielle et de l'apprentissage automatique.

Le supercalculateur porte le nom du mont Toubkal, point culminant du Maroc situé dans le Haut Atlas. Son déploiement constitue une étape importante dans le développement des capacités nationales et africaines de calcul haute performance (\textit{High Performance Computing}, HPC). L'infrastructure est destinée à soutenir les activités scientifiques nécessitant des ressources de calcul importantes.

\begin{figure}[htbp]
    \centering
    \includegraphics[width=0.7\textwidth]{Images_/supercalculateur_toubkal.png}
    \caption{Supercalculateur Toubkal de l'Université Mohammed VI Polytechnique}
    \label{fig:toubkal}
\end{figure}

Le supercalculateur Toubkal illustre la convergence croissante entre le calcul haute performance et l'intelligence artificielle. Les architectures modernes d'apprentissage profond reposent largement sur des accélérateurs matériels, notamment des processeurs graphiques (\textit{Graphics Processing Units}, GPU), permettant d'exécuter efficacement un grand nombre d'opérations matricielles en parallèle. Cette capacité de calcul devient particulièrement importante dans le contexte des modèles de grande taille et des applications nécessitant le traitement de volumes importants de données.

Toubkal dispose d'une infrastructure de calcul et de stockage destinée à répondre à des besoins scientifiques variés. Sa capacité permet notamment d'exécuter des simulations numériques, de traiter de grandes quantités de données et de réaliser des expérimentations liées à l'apprentissage automatique et à l'intelligence artificielle.

L'existence d'une telle infrastructure constitue un atout important pour l'écosystème scientifique de l'UM6P. Elle permet de réduire la dépendance à l'égard de ressources de calcul externes et facilite l'expérimentation de méthodes nécessitant des ressources matérielles importantes.

\subsection{Missions et domaines d'expertise}

\subsubsection{Développement d'un écosystème national et international en intelligence artificielle}

La mission d'AI Movement dépasse le cadre strict de la recherche académique. Le centre cherche à structurer un écosystème capable de réunir les compétences nationales et internationales dans le domaine de l'intelligence artificielle et des sciences des données. Cette ambition passe notamment par l'attraction d'experts internationaux, la création de partenariats scientifiques et le développement de collaborations entre chercheurs marocains et chercheurs étrangers.

Cette stratégie permet de favoriser le transfert de connaissances, l'émergence de nouvelles compétences et la réalisation de projets de recherche multidisciplinaires. Elle contribue également à renforcer la position du Maroc comme acteur régional dans le domaine de l'intelligence artificielle.

Le centre développe ainsi des collaborations avec des universités, des organismes internationaux, des entreprises et des institutions publiques. Ces partenariats permettent d'inscrire les activités de recherche du centre dans des problématiques dépassant les frontières nationales et d'aborder l'intelligence artificielle sous différents angles : scientifique, industriel, économique, éducatif, environnemental et sociétal.

\subsubsection{Recherche scientifique et développement technologique}

L'une des principales missions d'AI Movement consiste à associer recherche théorique, recherche appliquée, développement technologique, innovation et transfert industriel. Le centre cherche ainsi à rapprocher les laboratoires de recherche des entreprises et des problématiques concrètes rencontrées par les acteurs économiques et sociaux.

Cette approche favorise le développement de solutions d'intelligence artificielle adaptées aux besoins réels et aux spécificités des différents secteurs. Elle permet également de transformer les résultats issus de la recherche scientifique en prototypes, outils logiciels, méthodes ou solutions pouvant être transférés vers le monde industriel ou utilisés au bénéfice de la société.

Les travaux réalisés au sein de l'écosystème AI Movement couvrent ainsi un éventail large de domaines. Parmi les principales thématiques figurent notamment l'apprentissage automatique, l'apprentissage profond, la vision par ordinateur, le traitement automatique du langage naturel, les systèmes multi-agents, les agents intelligents, l'intelligence artificielle hybride, les systèmes autonomes, l'IA explicable ainsi que différentes applications interdisciplinaires.

Cette diversité scientifique permet au centre d'aborder des problématiques nécessitant le développement de méthodes capables d'exploiter des données hétérogènes et de répondre à des environnements complexes. Elle favorise également les interactions entre différentes disciplines de l'intelligence artificielle.

\subsubsection{Intelligence artificielle hybride et systèmes intelligents}

Une caractéristique importante de l'approche scientifique défendue par AI Movement réside dans la complémentarité entre différents paradigmes de l'intelligence artificielle. Le centre met notamment en avant une approche hybride et intégrative visant à construire des systèmes d'IA capables de tirer parti de plusieurs familles de méthodes et de connaissances.

Cette orientation est particulièrement pertinente dans le contexte actuel, où les systèmes d'intelligence artificielle doivent de plus en plus intégrer des données hétérogènes, combiner apprentissage et raisonnement, interagir avec leur environnement et prendre en compte des contraintes spécifiques liées aux domaines d'application.

L'intelligence artificielle hybride peut ainsi permettre d'associer, selon les besoins, des approches fondées sur l'apprentissage automatique et des approches symboliques ou à base de connaissances. Elle ouvre également la voie au développement de systèmes intelligents plus interprétables, plus adaptables et mieux intégrés aux environnements complexes.

Les systèmes multi-agents constituent également un domaine historiquement associé aux travaux de la Professeure Amal El Fallah Seghrouchni. Ces systèmes reposent sur la coordination de plusieurs agents autonomes capables de percevoir leur environnement, de communiquer, de coopérer ou de négocier afin d'atteindre des objectifs individuels ou collectifs.

\subsubsection{Formation et développement des compétences}

La formation constitue un autre axe important des missions d'AI Movement. Le centre vise à favoriser l'appropriation des technologies d'intelligence artificielle par différentes catégories de la population, depuis les étudiants et jeunes chercheurs jusqu'aux professionnels et aux acteurs économiques.

Cette mission répond à un enjeu stratégique majeur : le développement de l'intelligence artificielle dépend non seulement de la disponibilité d'infrastructures de calcul et de données, mais également de l'existence de compétences humaines capables de concevoir, développer, déployer et évaluer ces technologies.

Dans cette perspective, AI Movement participe à des activités de formation, d'encadrement scientifique, de séminaires, de conférences et d'accompagnement de doctorants et de jeunes chercheurs. Les présentations de travaux réalisées par les doctorants et les chercheurs du centre constituent notamment un moyen de favoriser le partage des connaissances et la diffusion des résultats scientifiques.

Le centre accorde également une attention à la formation continue et à la sensibilisation aux transformations induites par l'intelligence artificielle. Cette démarche vise à permettre à différents secteurs de la société de mieux comprendre les possibilités offertes par l'IA ainsi que les transformations organisationnelles et professionnelles susceptibles d'en découler.

\subsubsection{Transfert technologique, innovation et entrepreneuriat}

AI Movement cherche également à favoriser le transfert des résultats de la recherche vers le monde socio-économique. Cette orientation repose sur le développement de relations étroites avec les entreprises et sur la création de synergies entre recherche scientifique, innovation et besoins industriels.

L'objectif consiste notamment à transformer les résultats scientifiques en technologies et produits à forte valeur ajoutée, à favoriser la création de projets innovants et à encourager l'entrepreneuriat technologique. Le centre considère ainsi les start-ups et les jeunes entrepreneurs comme des acteurs importants de l'écosystème de l'intelligence artificielle.

Les collaborations avec des entreprises de différents secteurs illustrent cette volonté de rapprocher la recherche et les applications industrielles. Les domaines concernés peuvent notamment inclure l'énergie, l'industrie, la finance, l'agriculture, la santé et les services numériques.

\subsubsection{Applications sectorielles de l'intelligence artificielle}

L'approche d'AI Movement est également caractérisée par une volonté d'appliquer l'intelligence artificielle à des problématiques concrètes. Les domaines d'application peuvent concerner notamment l'agriculture, l'énergie, la santé, l'industrie, l'environnement, les infrastructures, l'éducation, la finance, la gouvernance et différents services numériques.

Cette orientation répond à une logique de recherche orientée vers l'impact. L'objectif n'est pas uniquement de développer de nouveaux algorithmes, mais également d'étudier comment ceux-ci peuvent contribuer à résoudre des problèmes spécifiques rencontrés par les sociétés et les organisations.

L'intelligence artificielle peut par exemple être utilisée pour améliorer la maintenance prédictive des infrastructures industrielles, analyser de grandes quantités de données, optimiser des processus, assister les opérateurs humains, détecter des anomalies, automatiser certaines tâches ou encore faciliter la prise de décision.

Cette dimension applicative constitue un élément essentiel de la stratégie du centre, puisqu'elle permet de rapprocher les avancées scientifiques des besoins réels des secteurs économiques et sociaux.

\subsubsection{Intelligence artificielle responsable et éthique}

Le développement de l'intelligence artificielle soulève parallèlement des questions scientifiques, économiques, juridiques et éthiques. La performance d'un système d'IA ne constitue donc pas l'unique critère permettant d'évaluer sa pertinence. Des considérations telles que la transparence, l'équité, la protection des données, la sécurité, la robustesse, l'explicabilité, l'inclusion et la responsabilité doivent également être prises en compte.

Cette dimension occupe une place importante dans la vision portée par AI Movement. Le centre considère que le développement de l'intelligence artificielle doit être accompagné d'une réflexion sur ses conséquences et sur les conditions permettant de garantir une utilisation responsable de ces technologies.

La question de l'inclusion constitue également un axe de réflexion important. Les systèmes d'intelligence artificielle sont susceptibles de reproduire ou d'amplifier certains biais présents dans les données utilisées pour leur conception. Il devient donc nécessaire de prendre en considération la diversité des populations, des langues, des contextes culturels et des environnements dans lesquels les systèmes sont déployés.

Dans le contexte africain, cette problématique revêt une importance particulière. Les modèles d'intelligence artificielle développés principalement à partir de données provenant d'autres régions du monde peuvent présenter des performances variables lorsqu'ils sont appliqués à des populations, langues ou contextes culturels différents. Le développement d'une expertise locale permet ainsi de mieux prendre en compte ces spécificités.

L'IA responsable implique également de porter une attention particulière à l'explicabilité des modèles, à la protection des données personnelles, à la sécurité des systèmes et à la responsabilité associée aux décisions automatisées. Ces considérations deviennent particulièrement importantes lorsque l'intelligence artificielle est utilisée dans des domaines sensibles tels que la santé, la finance, l'éducation ou les services publics.

\subsubsection{Une dimension africaine et internationale}

L'une des spécificités d'AI Movement réside enfin dans son ambition de contribuer au développement de l'intelligence artificielle à l'échelle africaine. Le centre ne se limite pas à une mission nationale : il cherche à participer à la constitution d'un écosystème africain de compétences, de recherche, de formation et d'innovation dans le domaine de l'IA.

Cette ambition se traduit par des partenariats et des initiatives internationales visant notamment à renforcer les capacités africaines dans le domaine de l'intelligence artificielle et à favoriser l'utilisation de ces technologies au service du développement durable.

Les collaborations avec des organisations internationales et des institutions de recherche permettent au centre de participer à des réseaux scientifiques dépassant le cadre marocain. Ces échanges contribuent à la circulation des connaissances, au développement de projets communs et à l'intégration des chercheurs et étudiants marocains dans des communautés scientifiques internationales.

L'ensemble de ces dimensions --- recherche, formation, innovation, transfert technologique, infrastructures de calcul, partenariats internationaux et réflexion éthique --- contribue à faire d'AI Movement un acteur structurant de l'écosystème marocain de l'intelligence artificielle.

Le centre s'inscrit ainsi dans une vision où l'intelligence artificielle constitue simultanément un domaine scientifique de pointe, un levier d'innovation technologique et un outil susceptible de contribuer au développement économique et social du Maroc et du continent africain.

\newpage

\section{Contexte du projet et problématique}

\subsection{L'essor des Modèles Vision-Langage (VLM) et le défi de l'échelle}

L'intelligence artificielle générative traverse actuellement une phase de transition majeure, marquée par le passage des modèles de langage purement textuels (LLM) vers des architectures multimodales : les Modèles Vision-Langage (VLM). Ces modèles fondationnels sont capables d'ingérer, de comprendre et de raisonner simultanément sur des flux de textes et d'images. Cependant, cette capacité cognitive accrue s'accompagne d'une exigence matérielle sans précédent.

L'entraînement de ces modèles à l'état de l'art nécessite le traitement de jeux de données massifs atteignant plusieurs pétaoctets, composés de milliards de paires image-texte. À cette échelle, l'infrastructure de calcul (telle que le supercalculateur Toubkal) déploie des grappes de processeurs graphiques (GPU) extrêmement performants. Néanmoins, une problématique fondamentale émerge : l'évolution de la puissance de calcul des accélérateurs matériels (GPU) a largement surpassé l'évolution de la bande passante des sous-systèmes de stockage et de la mémoire centrale (RAM).

\subsection{Le goulot d'étranglement des données (Data Loading Bottleneck)}

Dans un pipeline d'entraînement standard, le processeur central (CPU) est chargé de lire les données brutes depuis le disque, de décoder les images (souvent compressées aux formats JPEG ou PNG), de tokeniser le texte, d'appliquer des augmentations ou des normalisations, puis d'assembler ces éléments en lots (\textit{batches}) avant de les transférer vers la mémoire du GPU via le bus PCI-Express.

Historiquement, des solutions comme le \texttt{DataLoader} natif du framework PyTorch ont été conçues pour gérer ce flux. Or, ces solutions rencontrent aujourd'hui des limites architecturales sévères lorsqu'elles sont confrontées à des données multimodales à haute fréquence :
\begin{itemize}
    \item \textbf{Verrouillage global et parallélisme limité :} En Python, le \textit{Global Interpreter Lock} (GIL) empêche un véritable parallélisme multi-thread. Le recours au \textit{multiprocessing} engendre une surcharge mémoire (\textit{overhead}) considérable due à la duplication des processus et à la sérialisation des objets inter-processus.
    \item \textbf{Copies mémoire redondantes :} L'utilisation des appels système standards (comme \texttt{read()}) implique de multiples copies des données entre l'espace noyau (\textit{kernel space}), l'espace utilisateur (\textit{user space}) et la mémoire paginée, saturant ainsi la bande passante de la RAM avant même que les données n'atteignent le GPU.
\end{itemize}

\subsection{La problématique centrale : La "Famine du GPU" (GPU Starvation)}

La conséquence directe de ce goulot d'étranglement est un phénomène critique connu sous le nom de \textbf{GPU Starvation} (famine du GPU). Le processeur graphique, capable de réaliser des milliers de milliards d'opérations en virgule flottante par seconde (TFLOPs), termine le calcul de l'étape d'optimisation (\textit{forward} et \textit{backward pass}) bien avant que le CPU n'ait préparé le \textit{batch} suivant. Le GPU se retrouve alors dans un état d'inactivité (\textit{idle}), attendant les données. Ce gaspillage de ressources matérielles coûteuses réduit drastiquement l'efficacité (ou \textit{Hardware Flops Utilization}) des clusters de calcul.

À ce défi de l'ingestion des données s'ajoute une inefficacité algorithmique : le \textbf{Static Padding} (remplissage statique). Pour que le GPU puisse traiter des lots de données sous forme matricielle homogène, les approches traditionnelles forcent le redimensionnement de toutes les images à une résolution fixe maximale (ex: $384 \times 384$) et complètent les textes courts par des \textit{tokens} vides jusqu'à une longueur maximale (ex: 512). Sachant que le mécanisme d'attention au cœur des architectures Transformers possède une complexité temporelle et spatiale quadratique $\mathcal{O}(N^2)$, le calcul sur ces données artificiellement gonflées par le padding entraîne un gaspillage massif de FLOPs et de mémoire VRAM.

\subsection{Vision et objectifs du projet}

Face à ces limites technologiques et algorithmiques, ce projet de fin d'études se propose de repenser intégralement l'architecture d'ingestion des données pour l'entraînement des modèles Vision-Langage. 

L'objectif principal est de concevoir et d'implémenter un \textbf{pipeline de chargement distribué à très haute performance (HPC)}, capable de supprimer la famine du GPU. Pour y parvenir, le projet s'articule autour des axes suivants :
\begin{enumerate}
    \item \textbf{Ingénierie logicielle bas niveau :} Remplacer les lectures I/O classiques par une implémentation C++ asynchrone exploitant la projection mémoire (\texttt{mmap}) et le principe du \textit{Zero-Copy} pour éviter la saturation de la RAM.
    \item \textbf{Optimisation algorithmique :} Éliminer le gaspillage de calcul via une stratégie de \textit{Dynamic Padding}, préservant le ratio d'aspect natif des images et réduisant drastiquement le nombre d'opérations quadratiques dans le GPU.
    \item \textbf{Scalabilité et Distribution :} Concevoir un format binaire personnalisé (\textit{Shards}) et un mécanisme de partitionnement déterministe garantissant un entraînement distribué (DDP) efficace et robuste à grande échelle.
\end{enumerate}

Cette problématique, à l'intersection de l'architecture des systèmes, du calcul haute performance et de l'optimisation des algorithmes d'apprentissage profond, appelle une solution hybride alliant la flexibilité de Python à la puissance brute du C++, dont la conception et la validation scientifique feront l'objet des chapitres suivants.

% PAGE DE CHAPITRE STYLISÉE AVEC LOGOS ET TITRE CENTRÉ
\begin{titlepage}
  \thispagestyle{empty} % Supprime les en-têtes/pieds pour cette page

  % LOGOS HAUT DE PAGE
  \begin{minipage}{0.48\textwidth}
  \includegraphics[width=7cm]{Images_/logo_ensam_mek.png}\\[0.2cm]
  \textbf{}
\end{minipage}
\hfill
\begin{minipage}{0.48\textwidth}
  \raggedleft
  % Pense à ajouter le logo de AI Movement dans ton dossier Images_
  \includegraphics[width=5.5cm]{Images_/aimovement_logo.png}\\[0.4cm]
  \textbf{} \\
\end{minipage}

  % ESPACE VIDE POUR PLACEMENT VERTICAL AU CENTRE
  \vspace*{5cm}

  % TITRE AU CENTRE
  \begin{center}
    \rule{\linewidth}{0.5pt}\vspace{0.7cm}

    {\Huge\bfseries\color{blue}
    Chapitre 2 :\\[0.6em]
    Cadre Théorique et Fondements Mathématiques du Chargement de Données HPC}

    \vspace{0.7cm}
    \rule{\linewidth}{0.5pt}
  \end{center}

  \vfill

\end{titlepage}

\chapter{Cadre Théorique et Fondements Mathématiques du Chargement de Données HPC}
\label{chap:cadre_theorique}

\section{Le goulot d'étranglement des données dans l'entraînement des VLM}
\label{sec:goulot}

L'entraînement des modèles Vision-Langage (VLM) repose sur une hypothèse implicite trop souvent négligée dans la littérature : celle d'un flux de données parfaitement synchrone avec le calcul du GPU. En pratique, cette hypothèse s'effondre dès que la complexité du pipeline de préparation des données dépasse la capacité du GPU à consommer les échantillons au même rythme. Ce chapitre pose les fondements théoriques qui justifient, étape par étape, les choix d'architecture développés dans la suite de ce mémoire.

Un batch d'entraînement pour un VLM traverse une chaîne de transformations avant d'atteindre le GPU : décodage de l'image compressée (JPEG, PNG), redimensionnement et normalisation, tokenisation du texte associé, collation en un tenseur de batch homogène, puis transfert hôte-vers-périphérique (\textit{Host-to-Device}, H2D) via le bus PCIe. Chacune de ces étapes est exécutée sur CPU, à l'exception du transfert final, et chacune constitue un point de contention potentiel. Le phénomène de \textbf{GPU starvation} (famine du GPU) survient lorsque le GPU termine le calcul d'un pas d'entraînement avant que le pipeline CPU n'ait produit le batch suivant : le GPU reste alors inactif, et le débit effectif de l'entraînement (mesuré en échantillons/seconde) est plafonné non plus par la puissance de calcul disponible, mais par la vitesse du pipeline de données.

\subsection{Le modèle Roofline appliqué au pipeline de données}

Le modèle Roofline, introduit initialement pour caractériser la performance des noyaux de calcul, fournit un cadre pertinent pour formaliser ce déséquilibre. Il postule que le débit atteignable d'une opération est borné par le minimum entre la performance de calcul crête du processeur et le produit de la bande passante mémoire par l'intensité arithmétique de l'opération :

\begin{equation}
\label{eq:roofline}
P_{\text{atteignable}} = \min\bigl(P_{\text{pic}},\; BW \times I\bigr)
\end{equation}

où $P_{\text{pic}}$ désigne la performance de calcul crête (FLOPs/s), $BW$ la bande passante mémoire (octets/s), et $I$ l'intensité arithmétique, définie comme le rapport entre le nombre d'opérations flottantes exécutées et le nombre d'octets transférés depuis la mémoire ($I = \text{FLOPs}/\text{octets}$).

Un VLM combinant un encodeur visuel de type Vision Transformer (ViT) et un modèle de langage (LLM) présente une intensité arithmétique élevée : les opérations d'attention et les projections linéaires réutilisent massivement les poids chargés en mémoire, ce qui place ce type de charge dans le régime \textit{compute-bound} du diagramme Roofline, à droite du point d'inflexion (\textit{ridge point}). À l'inverse, le pipeline de préparation des données décodage d'image, redimensionnement, écriture/lecture disque présente une intensité arithmétique faible : peu d'opérations sont effectuées par octet transféré depuis le stockage. Ce pipeline est donc structurellement \textit{I/O-bound}. La coexistence de ces deux régimes sur la même machine crée une asymétrie fondamentale : le sous-système qui limite le débit global n'est presque jamais celui qui exécute le plus de calcul, mais celui qui déplace le plus d'octets par unité de temps utile.

\begin{figure}[htbp]
  \centering
  \begin{tikzpicture}[scale=1]
    \begin{axis}[
      xmode=log, ymode=log,
      xlabel={Intensité arithmétique $I$ (FLOPs/octet)},
      ylabel={Performance (FLOPs/s)},
      xmin=0.01, xmax=1000,
      ymin=1, ymax=1000,
      legend pos=south east,
      width=0.85\textwidth, height=6cm,
      grid=major
    ]
    \addplot[domain=0.01:10, samples=100, thick, blue] {50*x};
    \addlegendentry{Borne mémoire ($BW \times I$)}
    \addplot[domain=10:1000, samples=2, thick, red] {500};
    \addlegendentry{Borne calcul ($P_{\text{pic}}$)}
    \addplot[only marks, mark=*, mark size=2.5pt, color=orange] coordinates {(0.15,7.5)};
    \addlegendentry{Pipeline de données (I/O-bound)}
    \addplot[only marks, mark=square*, mark size=2.5pt, color=teal] coordinates {(50,500)};
    \addlegendentry{Bloc ViT/LLM (compute-bound)}
    \end{axis}
  \end{tikzpicture}
  \caption{Positionnement schématique du pipeline de données et du calcul VLM sur un diagramme Roofline.}
  \label{fig:roofline}
\end{figure}

\subsection{Condition formelle de non-famine}

Soit $\tau_{\text{step}}$ le temps de calcul d'un pas d'entraînement (\textit{forward} + \textit{backward} + mise à jour des poids) et $\tau_{\text{batch\_load}}$ le temps nécessaire au pipeline CPU pour produire un batch complet. La condition nécessaire et suffisante pour éviter la famine du GPU s'écrit :

\begin{equation}
\label{eq:non_famine}
\tau_{\text{batch\_load}} < \tau_{\text{step\_compute}}
\end{equation}

Cette condition peut être reformulée en termes de débit. En notant $b$ la taille du batch et $r$ le nombre de pas par seconde soutenu par le GPU, le taux de consommation de données du système s'exprime :

\begin{equation}
\label{eq:consommation}
\tau_{\text{consommation}} = b \times r \quad \text{(échantillons/seconde)}
\end{equation}

Le chargeur de données doit fournir un débit soutenu strictement supérieur à $\tau_{\text{consommation}}$ pour ne jamais devenir le facteur limitant. Cette formalisation motive directement la démarche méthodologique adoptée dans ce travail : mesurer $\tau_{\text{batch\_load}}$ de manière comparative entre notre implémentation et une solution de référence de l'écosystème, sur des corpus réels de type VQA, afin que les mesures reflètent une distribution réaliste de tailles d'image et de longueurs de texte plutôt qu'un jeu de données synthétique artificiellement favorable. Le protocole expérimental correspondant et ses résultats sont détaillés au chapitre consacré à l'évaluation.

\subsection{Théorie des files d'attente et dimensionnement du prefetch}

Le couple chargeur/GPU peut être modélisé comme un système producteur-consommateur au sens de la théorie des files d'attente : le chargeur agit comme serveur produisant des batches, le GPU comme client les consommant. En présence de variance dans le temps de production (\textit{jitter}), une file de prefetch de profondeur nulle expose immédiatement le système à la famine dès qu'un seul batch dépasse le temps moyen. La profondeur de file nécessaire pour absorber cette variance sans provoquer d'attente est approximée par :

\begin{equation}
\label{eq:prefetch_depth}
d_{\text{prefetch}} \approx \frac{\sigma_{\text{latence}}}{\tau_{\text{step}}}
\end{equation}

où $\sigma_{\text{latence}}$ est l'écart-type du temps de production d'un batch. Cette relation établit qu'un pipeline à latence très variable (par exemple à cause de tailles d'image hétérogènes ou de contention disque) nécessite un tampon de prefetch profond, au prix d'une empreinte mémoire accrue.

Ce dimensionnement peut être rattaché formellement à la \textbf{loi de Little}, résultat fondamental de la théorie des files d'attente valable pour tout système stationnaire : le nombre moyen d'éléments présents dans le système, $L$, est égal au produit du taux d'arrivée moyen $\lambda$ par le temps de séjour moyen $W_{\text{séjour}}$ d'un élément dans le système,

\begin{equation}
\label{eq:little}
L = \lambda \times W_{\text{séjour}}
\end{equation}

Appliquée au tampon de prefetch, $\lambda$ correspond au taux de production de batches par le pipeline CPU et $W_{\text{séjour}}$ au temps moyen qu'un batch passe dans le tampon avant consommation par le GPU. La loi de Little justifie qu'à taux de production fixé, augmenter la profondeur de la file (donc $L$) revient mécaniquement à augmenter la latence de bout en bout $W_{\text{séjour}}$ subie par chaque batch un compromis direct entre robustesse au jitter (via \eqref{eq:prefetch_depth}) et fraîcheur des données consommées, particulièrement pertinent lorsque le pipeline effectue un mélange (\textit{shuffling}) en flux.

\section{Hiérarchie mémoire, appels système et zero-copy : la théorie de la projection mémoire}
\label{sec:mmap}

La deuxième source de surcoût identifiée dans la chaîne de préparation des données concerne l'accès au stockage lui-même. Deux stratégies s'opposent : la lecture classique par appel système et la projection mémoire (\textit{memory mapping}) du fichier.

\subsection{Coût de la lecture classique}

Un appel système de lecture (\texttt{read}) déclenche une transition de mode processeur du niveau utilisateur (\textit{ring 3}) vers le niveau noyau (\textit{ring 0}), coûteuse en cycles processeur du fait de la sauvegarde de contexte et du changement de table de pages. Le noyau copie ensuite les données depuis le disque vers le \textit{page cache} (première copie mémoire), puis du \textit{page cache} vers le tampon utilisateur fourni par l'application (seconde copie mémoire). Le coût total par lecture comprend donc un surcoût fixe de syscall et un coût variable proportionnel à la taille lue :

\begin{equation}
\label{eq:read_cost}
T_{\text{read}}(N) \approx \frac{N}{BW_{\text{disque}}} + \frac{2N}{BW_{\text{mem}}} + \frac{N}{\text{taille\_lecture}} \times t_{\text{syscall}}
\end{equation}

Le terme $N/BW_{\text{disque}}$ correspond au transfert physique, $2N/BW_{\text{mem}}$ aux deux copies mémoire successives, et le dernier terme au coût amorti des transitions ring 3 → ring 0, multiplié par le nombre d'appels système nécessaires pour lire $N$ octets.

\subsection{Zéro-copie via projection mémoire}

La projection mémoire d'un fichier (\texttt{mmap}) adopte une stratégie radicalement différente : plutôt que de copier les données, le noyau projette directement les pages du fichier dans l'espace d'adressage virtuel du processus appelant. Aucune copie n'est effectuée au moment de l'appel celui-ci se contente d'établir une correspondance entre adresses virtuelles et pages physiques du \textit{page cache}. L'accès effectif aux données ne survient que lors du premier déréférencement, qui déclenche un \textbf{défaut de page} (\textit{page fault}) : le processeur consulte la table des pages TLB, constate l'absence de correspondance physique valide, et délègue au noyau la résolution du défaut chargement de la page depuis le disque vers le \textit{page cache} si elle n'y est pas déjà, puis mise à jour de la table de pages. Ce mécanisme constitue un chargement à la demande (\textit{lazy loading}) : seules les pages effectivement accédées sont matérialisées en mémoire physique.

Le coût amorti de cette stratégie s'écrit :

\begin{equation}
\label{eq:mmap_cost}
T_{\text{mmap}}(N) \approx \frac{N}{BW_{\text{page\_cache}}} + \frac{N}{\text{taille\_page}} \times t_{\text{fault}}
\end{equation}

où le premier terme correspond au débit du \textit{page cache} (aucune copie applicative n'étant nécessaire, le CPU ne recopie jamais les données vers un tampon en \textit{heap}) et le second au coût amorti d'un défaut de page par page mémoire (typiquement 4 Ko) plutôt que par appel système. La comparaison de \eqref{eq:read_cost} et \eqref{eq:mmap_cost} montre que la projection mémoire élimine structurellement une copie mémoire complète et amortit le coût de transition noyau sur des unités de granularité plus grande (la page) plutôt que sur chaque appel applicatif.

Le comportement par défaut du noyau face à un défaut de page résolution synchrone page paSr page peut être optimisé de deux manières complémentaires. D'une part, une demande de peuplement anticipé de la table de pages, formulée au moment même de la projection, déclenche le \textit{readahead} de la totalité du fichier avant même le premier accès applicatif ce qui déplace le coût des défauts de page du chemin critique de lecture vers l'ouverture du fichier, où il peut être recouvert par d'autres opérations. D'autre part, un conseil d'accès (\textit{advice}) donné au noyau après la projection, du type « accès séquentiel », informe l'algorithme de lecture anticipée que les accès suivront l'ordre des octets du fichier, ce qui ajuste la politique interne de \textit{readahead} et l'agressivité de l'éviction du cache en conséquence.

\begin{table}[htbp]
  \centering
  \renewcommand{\arraystretch}{1.5}
  \begin{tabular}{|p{4.5cm}|p{5cm}|p{5cm}|}
    \hline
    \textbf{Critère} & \textbf{Lecture classique (\texttt{read})} & \textbf{Projection mémoire (\texttt{mmap})} \\
    \hline
    Copies mémoire par accès & 2 (disque $\rightarrow$ cache, cache $\rightarrow$ tampon) & 0 (accès direct au \textit{page cache}) \\
    \hline
    Granularité du surcoût noyau & par appel système & par page (typiquement 4 Ko) \\
    \hline
    Déclenchement du transfert & explicite (appel bloquant) & implicite (défaut de page) \\
    \hline
    Chargement & anticipé, contrôlé par l'appelant & paresseux (\textit{lazy}), à la demande \\
    \hline
    Partage inter-processus du cache & indirect (via le \textit{page cache}) & direct (pages physiques partagées) \\
    \hline
  \end{tabular}
  \vspace{0.2cm}
  \caption{Comparaison théorique des deux stratégies d'accès au fichier.}
  \label{tab:read_vs_mmap}
\end{table}

\subsection{Alignement et localité spatiale}

Le gain de la projection mémoire n'est pleinement exploité que si les accès respectent la granularité des lignes de cache du processeur (64 octets sur la quasi-totalité des architectures x86-64 et ARM64 contemporaines). Un échantillon dont l'adresse de début n'est pas un multiple de la taille d'une ligne de cache chevauche potentiellement deux lignes distinctes : un unique accès logique déclenche alors deux transactions mémoire physiques (phénomène de \textit{cache-line split}), dégradant le débit effectif indépendamment du mécanisme d'accès retenu. C'est cette contrainte microarchitecturale qui justifie, dans notre format de sérialisation binaire, l'alignement systématique de chaque échantillon sur une frontière de 64 octets : cette contrainte de conception, au niveau du format de fichier, est directement dérivée de la microarchitecture du processeur cible plutôt que d'un choix arbitraire.

Notre implémentation du lecteur exploite cette stratégie via une abstraction de fichier projeté en mémoire, en lecture seule et en mode de projection privée (garantissant une vue \textit{copy-on-write} propre au processus, sans propagation des écritures vers le fichier sous-jacent cohérent avec un rôle strictement consommateur de données déjà écrites), combinée aux mécanismes de peuplement anticipé et de conseil d'accès séquentiel décrits ci-dessus.

\section{Complexité quadratique du Transformer et théorie du padding}
\label{sec:padding}

\subsection{Coût algorithmique de l'attention}

Chaque couche d'un Transformer, qu'il s'agisse du ViT encodant l'image ou du LLM traitant le texte, exécute deux familles d'opérations dominantes. Les projections linéaires (requêtes, clés, valeurs, sortie, ainsi que les couches \textit{feed-forward}) ont un coût en $O(N \cdot d^2)$ FLOPs, où $N$ est la longueur de séquence et $d$ la dimension du modèle. Le calcul de l'attention proprement dit produit $QK^\top$ suivi de la pondération des valeurs $AV$ a un coût en $O(N^2 \cdot d)$ FLOPs, et la matrice d'attention intermédiaire occupe $O(N^2)$ en mémoire. Le coût total par couche s'écrit approximativement :

\begin{equation}
\label{eq:transformer_flops}
\text{FLOPs}_{\text{couche}} \approx 2N d^2 + 4N^2 d
\end{equation}

Pour un ViT, $N$ correspond au nombre de patches issus du découpage de l'image, $N = (H/P) \times (W/P)$, où $H$ et $W$ sont les dimensions de l'image et $P$ la taille du patch. Pour un LLM, $N$ est directement la longueur de la séquence de tokens.

\subsection{Le gaspillage du padding statique}

Une stratégie de padding statique fixe $N$ à une valeur maximale unique pour l'ensemble du corpus par exemple $N_{\text{fixe}} = (384/16)^2 = 576$ patches pour toutes les images d'un ViT opérant sur des entrées de $384 \times 384$ pixels avec des patches de $16 \times 16$, et $\text{seq\_len} = 512$ pour tous les textes. Or la distribution réelle des tailles d'image et des longueurs de texte d'un corpus VQA est très majoritairement concentrée en deçà de ce maximum. Étant donné la dépendance quadratique du coût de l'attention en $N$, chaque échantillon dont la taille effective est inférieure au maximum fixe gaspille une fraction du calcul qui croît quadratiquement avec l'écart entre taille effective et taille maximale.

Le padding dynamique, appliqué au niveau de la collation d'un batch plutôt qu'au niveau du corpus entier, ramène $N$ au maximum effectivement présent \emph{dans le batch courant} plutôt qu'au maximum du corpus. Le gain relatif de calcul entre les deux stratégies suit directement le rapport quadratique :

\begin{equation}
\label{eq:gain_padding}
\frac{\text{FLOPs}_{\text{statique}}}{\text{FLOPs}_{\text{dynamique}}} \approx \left(\frac{N_{\text{fixe}}}{N_{\text{batch\_max}}}\right)^2
\end{equation}

pour la composante attention, dominante lorsque $N$ est grand devant $d$. C'est cette relation qui motive, dans notre architecture, la collation dynamique au maximum du batch, ainsi que le redimensionnement préservant le ratio d'aspect en amont du pipeline — une stratégie qui limite la dispersion des tailles d'image plutôt que de la corriger a posteriori par un padding uniforme excessif.

\subsection{Réduction du volume I/O par stockage en entiers non signés}

Indépendamment du coût de calcul, le volume de données transférées depuis le stockage est directement proportionnel à la taille de l'élément de stockage utilisé. Une image stockée en flottant 32 bits occupe $3 \cdot H \cdot W \cdot 4$ octets (trois canaux, quatre octets par élément), contre $3 \cdot h \cdot w \cdot 1$ octet en entier non signé 8 bits (\texttt{uint8}). Le facteur de réduction I/O entre les deux représentations, indépendamment de la résolution effective, est exactement :

\begin{equation}
\label{eq:reduction_io}
\frac{\text{octets}_{\text{f32}}}{\text{octets}_{\text{uint8}}} = \frac{3 H W \times 4}{3 H W \times 1} = 4
\end{equation}

soit une réduction par un facteur 4 du volume à lire depuis le disque. Ce gain n'est exploitable sans perte de précision que si la normalisation classiquement $(x/255 - \mu)/\sigma$ est différée de l'écriture du fichier vers l'étape de collation, au moment où le tenseur est effectivement converti en flottant pour être consommé par le modèle. Ce report est théoriquement neutre du point de vue numérique, la multiplication scalaire étant associative : appliquer la normalisation avant ou après le stockage intermédiaire en entiers non signés produit un résultat final identique à la précision de calcul près, dès lors que ce stockage n'introduit lui-même aucune perte par rapport à la représentation d'origine des pixels (déjà quantifiée sur 8 bits par canal dans la quasi-totalité des formats image usuels). Cette différée de normalisation, couplée au stockage des dimensions réelles de chaque échantillon plutôt que d'une dimension fixe pour l'ensemble d'un fichier, est la condition nécessaire à l'exploitation effective du padding dynamique décrit ci-dessus.

\section{Parallélisme : lois de scalabilité, GIL et concurrence sans verrou}
\label{sec:parallelisme}

\subsection{Le GIL de CPython et la nécessité du multiprocessing}

L'interpréteur de référence CPython impose un verrou global (\textit{Global Interpreter Lock}, GIL) garantissant qu'un seul thread exécute du bytecode Python à un instant donné, quel que soit le nombre de cœurs physiques disponibles. Cette contrainte interdit tout gain de performance par parallélisme de threads pour du code CPU-bound écrit en Python pur le décodage d'image et la tokenisation, bien qu'apparemment parallélisables, ne bénéficient d'aucun \textit{speedup} réel sous un modèle à threads. La seule voie de parallélisation effective au niveau processus consiste à recourir à plusieurs processus système indépendants, chacun disposant de son propre interpréteur et donc de son propre GIL, la communication inter-processus s'effectuant via des files de messages ou de la mémoire partagée plutôt que par accès direct à des objets Python partagés.

Le \textit{speedup} observé en fonction du nombre de workers $n$ est borné par la loi d'Amdahl, qui tient compte de la fraction incompressible de travail séquentiel $(1-p)$ dans un pipeline de préparation de données, typiquement l'ingestion depuis la source et l'écriture finale des fichiers de sortie, opérations qui ne peuvent être parallélisées sans rupture d'ordre ou risque de corruption d'écriture concurrente :

\begin{equation}
\label{eq:amdahl}
S(n) = \frac{1}{(1-p) + \dfrac{p}{n}}
\end{equation}

L'efficacité par worker, $E(n) = S(n)/n$, décroît nécessairement avec $n$ dès lors que $p < 1$, ce qui borne le nombre de workers économiquement pertinent bien avant la saturation du nombre de cœurs physiques disponibles. La limite asymptotique de \eqref{eq:amdahl} quand $n \to \infty$ vaut $S_{\max} = 1/(1-p)$ : quelle que soit la puissance de calcul additionnelle investie, le \textit{speedup} ne peut dépasser l'inverse de la fraction séquentielle du pipeline ici, essentiellement le temps fixe d'ingestion et d'écriture séquentielle, difficilement parallélisable sans compromettre l'intégrité de l'ordonnancement des données produites. Ce résultat motive un choix de dimensionnement pratique : le nombre de workers optimal n'est pas le nombre de cœurs physiques disponibles, mais le plus petit $n$ au-delà duquel le gain marginal de \textit{speedup}, $S(n+1) - S(n)$, devient négligeable devant le coût marginal en pression mémoire et en contention du système de fichiers induit par des workers supplémentaires.

\subsection{Traitement par fenêtre bornée et mémoire constante}

Le découplage entre taille du corpus et empreinte mémoire du pipeline hors-ligne repose sur une stratégie de \textbf{fenêtrage} au niveau de l'orchestration : l'ensemble des $N$ échantillons est partitionné en $\lceil N / c \rceil$ blocs (\textit{chunks}) de taille bornée $c$, chaque bloc étant successivement soumis à un ensemble de workers réutilisés d'un bloc à l'autre (évitant le coût de réinitialisation répétée des ressources de traitement), puis écrit sur le support de stockage avant d'être explicitement libéré de la mémoire le bloc suivant n'étant chargé qu'à ce moment. À tout instant, seules les données d'au plus un bloc de taille $c$ sont matérialisées en mémoire simultanément, quelle que soit la taille totale $N$ du corpus :

\begin{equation}
\label{eq:memoire_bornee}
M_{\text{peak}} = O(c)
\end{equation}

indépendamment de $N$. Cette propriété est le principe qui distingue un pipeline en flux (\textit{streaming}) d'un pipeline qui chargerait l'intégralité du corpus en mémoire avant traitement condition indispensable pour le passage à l'échelle sur des corpus de plusieurs millions d'échantillons et elle fait l'objet, dans notre validation expérimentale, d'une vérification empirique consistant à exécuter le pipeline sur des corpus de tailles croissantes et à contrôler que le pic de mémoire tracée ne croît pas proportionnellement.

\subsection{Files SPSC sans verrou}

Pour la communication à haute fréquence entre un flux de préchargement asynchrone et le thread de consommation, l'usage d'un verrou classique introduirait, à la fréquence d'un \textit{push}/\textit{pop} par batch, un surcoût de synchronisation disproportionné par rapport au coût de l'opération protégée. Notre implémentation asynchrone recourt à une file \textbf{SPSC} (\textit{Single Producer Single Consumer}) sans verrou, structurée comme un tampon circulaire (\textit{ring buffer}) de capacité fixe $C$. Une capacité utile de $C-1$ éléments est effectivement exploitable : un emplacement sert de sentinelle pour distinguer sans ambiguïté l'état plein de l'état vide sans recourir à un compteur d'éléments séparé.

La structure repose sur deux indices atomiques, \texttt{head} (lu et incrémenté par le consommateur) et \texttt{tail} (lu et incrémenté par le producteur), tous deux progressant modulo $C$. Les conditions de plénitude et de vacuité de la file s'expriment :

\begin{align}
\text{file pleine} &\iff (\text{tail} + 1) \bmod C = \text{head} \label{eq:full}\\
\text{file vide} &\iff \text{tail} = \text{head} \label{eq:empty}
\end{align}

où \eqref{eq:full} est la condition testée avant toute tentative d'écriture, et \eqref{eq:empty} la condition testée avant toute tentative de lecture les deux opérations étant non bloquantes (\textit{wait-free}) : en cas d'échec, l'appelant reçoit un signal d'échec immédiat plutôt que d'attendre, à charge pour lui d'implémenter sa propre politique d'attente (\textit{spin} ou \textit{yield}).

La correction de cette structure sans verrou repose sur le modèle mémoire du C++ moderne (les types atomiques standard) et plus précisément sur la garantie d'ordonnancement \textit{acquire-release} : le producteur publie un nouvel élément en écrivant l'indice \texttt{tail} avec une sémantique de \textit{release}, garantissant que toutes les écritures antérieures dans l'élément lui-même sont visibles par tout thread effectuant une lecture correspondante avec une sémantique d'\textit{acquire} sur ce même indice établissant une relation \textit{happens-before} entre la production et la consommation de l'élément, sans nécessiter de verrou explicite.

Un second aspect microarchitectural est la protection contre le \textbf{faux partage} (\textit{false sharing}) : si les indices \texttt{head} et \texttt{tail} résident sur la même ligne de cache (64 octets), chaque écriture de l'un par son thread propriétaire invalide, au titre du protocole de cohérence de cache MESI, la copie de cette ligne détenue par le cœur exécutant l'autre thread y compris pour la partie de la ligne non concernée par l'écriture. Ce phénomène dégrade sévèrement le débit de la file en pratique, en dépit de son coût algorithmique théoriquement nul. Le padding explicite des deux indices sur des lignes de cache distinctes élimine cette invalidation croisée.

Cette structure générique instancie concrètement, dans notre chargeur asynchrone, la file de prefetch dont la profondeur a été introduite théoriquement en \eqref{eq:prefetch_depth} : un paramètre configurable désigne le nombre de batches maintenus prêts en mémoire pendant que le GPU consomme le batch courant un cas particulier de tamponnage multiple (\textit{multiple buffering}) directement dérivé du modèle producteur-consommateur de la section \ref{sec:goulot}.

\begin{figure}[htbp]
  \centering
  \begin{tikzpicture}[scale=1.2, transform shape]
    
    \foreach \i in {0,...,7} {
      \node[draw, thick, minimum width=1.2cm, minimum height=1.2cm] (c\i) at (\i*1.2,0) {};
    }
    
    \node[below=0.2cm] at (c1.south) {\textbf{head}};
    \node[below=0.2cm] at (c5.south) {\textbf{tail}};
    
    \draw[->, very thick, blue] (c1.north) -- ++(0,0.8) node[above, text=black] {consommateur lit};
    \draw[<-, very thick, red] (c5.north) -- ++(0,0.8) node[above, text=black] {producteur écrit};
    
  \end{tikzpicture}
  \vspace{0.3cm}
  \caption{Tampon circulaire SPSC : indices \texttt{head}/\texttt{tail} atomiques sur lignes de cache distinctes.}
  \label{fig:spsc}
\end{figure}

\section{Entraînement distribué et partitionnement déterministe des données}
\label{sec:distribue}

\subsection{Exigences formelles du partitionnement}

Dans un entraînement distribué DDP impliquant $W$ processus (ou rangs), chaque rang doit consommer, à chaque époque, une partition disjointe de l'ensemble des données. Quatre propriétés formelles doivent être garanties par le mécanisme de partitionnement, résumées au tableau \ref{tab:ddp_requirements}.

\begin{table}[htbp]
  \centering
  \renewcommand{\arraystretch}{1.5} 
  \begin{tabular}{|p{3.5cm}|p{5cm}|p{6.5cm}|}
    \hline
    \textbf{Propriété} & \textbf{Définition formelle} & \textbf{Conséquence d'une violation} \\
    \hline
    Couverture totale & $\bigcup_i A_i = S$ & Certains échantillons ne sont jamais vus par aucun rang \\
    \hline
    Absence de recouvrement & $A_i \cap A_j = \emptyset,\ i \neq j$ & Un échantillon est compté plusieurs fois par époque, biaisant l'estimation du gradient \\
    \hline
    Déterminisme & même graine, même époque $\Rightarrow$ même partition & Non-reproductibilité, reprise incohérente après un point de contrôle \\
    \hline
    Équilibrage & déséquilibre de taille entre rangs $\leq$ une unité de partition & Un rang devient systématiquement le goulot d'étranglement de la synchronisation collective de fin d'époque \\
    \hline
  \end{tabular}
  \vspace{0.2cm}
  \caption{Exigences formelles du partitionnement pour l'entraînement distribué (DDP).}
  \label{tab:ddp_requirements}
\end{table}

\subsection{Attribution strided et preuve de partition}

Le partitionnement est ici opéré non pas au niveau de l'échantillon individuel, mais au niveau du \textbf{fragment} (\textit{shard}) de données : chaque fragment, unité atomique de lecture et de vérification d'intégrité, est traité comme un élément indivisible de l'attribution. Le mécanisme retenu, dit \textit{strided} (répartition cyclique), assigne le fragment d'indice $j$ au rang $i$ selon la relation de congruence :

\begin{equation}
\label{eq:strided}
A_i = \{\, \sigma_j : j \equiv i \pmod{W} \,\}
\end{equation}

où $\sigma_j$ désigne le $j$-ième fragment dans un ordre fixe et déterministe. La preuve que cette attribution constitue effectivement une partition de l'ensemble des fragments découle directement des propriétés de la relation de congruence modulo $W$ : pour tout indice $j$, il existe un unique reste $i = j \bmod W \in \{0, \ldots, W-1\}$, ce qui garantit simultanément la couverture totale (tout fragment est affecté à exactement un rang) et l'absence de recouvrement (les classes de congruence modulo $W$ sont deux à deux disjointes par définition de la division euclidienne). Le déséquilibre \emph{en nombre de fragments} entre rangs est borné par construction :

\begin{equation}
\label{eq:desequilibre}
\left\lceil \frac{K}{W} \right\rceil - \left\lfloor \frac{K}{W} \right\rfloor \leq 1
\end{equation}

où $K$ est le nombre total de fragments. Il convient de noter que cette borne porte sur le \emph{nombre de fragments} attribués à chaque rang, et non directement sur le nombre d'échantillons : l'équilibrage en échantillons n'est garanti que dans la mesure où les fragments sont eux-mêmes de taille comparable, ce que la politique d'écriture par blocs de taille fixe (section \ref{sec:parallelisme}) tend à assurer en pratique. Un cas limite mérite également d'être formalisé : lorsque $W > K$ (plus de rangs que de fragments), la relation \eqref{eq:strided} produit $A_i = \emptyset$ pour $K \leq i < W$ certains rangs ne reçoivent alors aucun fragment et doivent être gérés explicitement côté boucle d'entraînement (par exemple en consommant un batch factice) pour ne pas bloquer la synchronisation collective.

Ces propriétés déterminisme, couverture totale, absence de recouvrement, comportement aux cas limites et respect de la borne de déséquilibre \eqref{eq:desequilibre} sont garanties par construction dans notre module de partitionnement distribué, et vérifiées de manière exhaustive dans notre suite de tests automatisés dédiée à l'entraînement distribué.

\subsection{Le manifeste comme contrat de coordination}

Un manifeste de métadonnées, structuré de manière portable entre processus, joue le rôle de contrat explicite entre les processus d'écriture (côté préparation des données) et les processus de lecture distribués (côté entraînement) : il fixe la correspondance entre rangs et fragments, ainsi qu'une somme de contrôle (\textit{checksum}) CRC32 par fragment permettant de détecter toute corruption survenue entre l'écriture et la lecture — un risque non négligeable à l'échelle de corpus de plusieurs téraoctets répartis sur des systèmes de fichiers réseau.

Le CRC32 repose sur l'arithmétique des polynômes à coefficients dans le corps fini $\mathbb{F}_2 = \{0,1\}$ (addition et multiplication modulo 2, sans retenue). Un message de $M$ bits $m_{M-1} \ldots m_1 m_0$ est interprété comme un polynôme $m(x) = \sum_{k=0}^{M-1} m_k x^k$. Le CRC32 est défini comme le reste de la division polynomiale de $m(x) \cdot x^{32}$ (le message décalé de 32 bits) par un polynôme générateur fixe $g(x)$ de degré 32 :

\begin{equation}
\label{eq:crc_def}
\text{CRC}(m) = \bigl(m(x) \cdot x^{32}\bigr) \bmod g(x)
\end{equation}

La norme ISO 3309, qui fixe ce générateur $g(x)$ (sous sa forme réfléchie usuelle pour un calcul bit-first-LSB), définit un code de détection et non de correction d'erreurs : contrairement à un code correcteur, il ne permet pas de localiser ni de réparer l'erreur, seulement de signaler son existence. Sa propriété fondamentale découle directement de la structure de l'anneau des polynômes sur $\mathbb{F}_2$ : une corruption du message équivaut à l'ajout d'un polynôme d'erreur $e(x)$ non nul, et cette erreur échappe à la détection si et seulement si $e(x)$ est un multiple de $g(x)$. Le choix d'un polynôme générateur primitif de degré 32 garantit que \emph{tout} polynôme d'erreur non nul de degré strictement inférieur à 32 (donc toute rafale d'erreurs, \textit{burst error}, de longueur $\leq 32$ bits) n'est jamais un multiple de $g(x)$, et est donc détecté à coup sûr ; un choix de $g(x)$ possédant un facteur $(x+1)$ garantit en outre la détection de toute erreur affectant un nombre impair de bits, puisqu'un polynôme d'erreur à poids de Hamming impair ne peut s'annuler en $x=1$.

La probabilité qu'une corruption de données produise malgré tout un CRC32 identique à l'original condition nécessaire pour qu'une corruption passe totalement inaperçue est approximée, pour des erreurs indépendantes du polynôme, par :

\begin{equation}
\label{eq:crc_proba}
P_{\text{non-détection}} \approx 2^{-32}
\end{equation}

soit environ $2{,}3 \times 10^{-10}$, une probabilité jugée suffisamment faible pour un usage de vérification d'intégrité de fragments, sans toutefois offrir la garantie cryptographique qu'apporterait une fonction de hachage plus coûteuse (SHA-256 par exemple) compromis jugé raisonnable étant donné le volume de données à vérifier et la nature accidentelle, et non adversariale, de la corruption visée. Cette vérification est matérialisée, dans notre format de sérialisation, par un champ de contrôle placé en fin de chaque fragment.

\bigskip
\noindent Ce chapitre a établi le socle théorique justifiant les décisions d'architecture développées dans la suite de ce mémoire : le modèle Roofline et la condition de non-famine motivent l'existence même du projet et sa démarche d'évaluation comparative ; la théorie de la projection mémoire justifie le choix du zéro-copie et de l'alignement sur les lignes de cache ; la complexité quadratique de l'attention justifie le padding dynamique et le stockage en entiers non signés ; la loi d'Amdahl et les structures sans verrou justifient l'architecture multiprocessus et la file de prefetch sans verrou ; enfin, la théorie du partitionnement déterministe justifie la conception du manifeste et son mécanisme de vérification par CRC32. Le chapitre suivant s'appuiera sur ce socle pour détailler l'architecture logicielle concrète du système.

\begin{titlepage}
  \thispagestyle{empty} % Supprime les en-têtes/pieds pour cette page

  % LOGOS HAUT DE PAGE
  \begin{minipage}{0.48\textwidth}
  \includegraphics[width=7cm]{Images_/logo_ensam_mek.png}\\[0.2cm]
  \textbf{}
\end{minipage}
\hfill
\begin{minipage}{0.48\textwidth}
  \raggedleft
  % Pense à ajouter le logo de AI Movement dans ton dossier Images_
  \includegraphics[width=5.5cm]{Images_/aimovement_logo.png}\\[0.4cm]
  \textbf{} \\
\end{minipage}

  % ESPACE VIDE POUR PLACEMENT VERTICAL AU CENTRE
  \vspace*{5cm}

  % TITRE AU CENTRE
  \begin{center}
    \rule{\linewidth}{0.5pt}\vspace{0.7cm}

    {\Huge\bfseries\color{blue}
    Chapitre 3 :\\[0.6em]
    Conception et Modélisation de l’Architecture}

    \vspace{0.7cm}
    \rule{\linewidth}{0.5pt}
  \end{center}

  \vfill

\end{titlepage}

\chapter{Conception et Modélisation de l'Architecture}
\label{chap:architecture}

\section{État de l'art et positionnement des solutions de chargement}
\label{sec:etat_art}

La conception d'un chargeur de données pour l'entraînement de modèles Vision-Langage ne s'opère pas dans un vide technique : plusieurs familles de solutions existent déjà pour l'alimentation de pipelines d'apprentissage à grande échelle. Leur examen critique, selon quatre axes déterminants pour notre cas d'usage, permet de délimiter précisément l'espace de conception laissé vacant et de justifier objectivement les choix architecturaux du chapitre.

Les quatre axes retenus sont : le \textbf{support multimodal natif} (capacité à encoder conjointement image et texte dans une même unité de stockage, plutôt que deux flux synchronisés séparément) ; le mode d'\textbf{accès aux données} (accès aléatoire à un échantillon arbitraire en temps constant, contre lecture strictement séquentielle en flux) ; le \textbf{coût de copie mémoire} induit par le chemin de lecture (nombre de copies entre le support de stockage et le tenseur consommé par le modèle) ; et la \textbf{compatibilité native avec l'entraînement distribué} (partitionnement déterministe intégré, plutôt que délégué entièrement à une bibliothèque tierce).

\begin{table}[htbp]
  \centering
  \renewcommand{\arraystretch}{1.5} % التيساع بين السطورا
  % كبرنا العبارات ديال الخانات باش الكلمات الطوال ما يتقطعوش
  \begin{tabular}{|p{4.5cm}|p{2.6cm}|p{2.8cm}|p{3.4cm}|p{2.2cm}|}
    \hline
    \textbf{Approche} & \textbf{Multimodal natif} & \textbf{Accès} & \textbf{Copies mémoire} & \textbf{DDP intégré} \\
    \hline
    Archives séquentielles (type TAR/TFRecord) & Partiel & Streaming seul & Élevé (désérialisation) & Externe \\
    \hline
    Format colonnaire à cache local & Partiel & Aléatoire (via cache) & Moyen (matérialisation) & Externe \\
    \hline
    Pipeline compilé JIT, orienté vision & Non (image seule) & Aléatoire & Faible & Externe \\
    \hline
    Format versionné en couches (\textit{data lake}) & Oui & Aléatoire (indexé) & Moyen (couche réseau) & Externe \\
    \hline
    Chargeur générique de référence (écosystème PyTorch) & Oui (via code utilisateur) & Aléatoire ou streaming & Élevé (plusieurs copies Python) & Externe \\
    \hline
    \textbf{Approche retenue dans ce travail} & \textbf{Oui} & \textbf{Aléatoire (table d'offsets)} & \textbf{Nul (zéro-copie)} & \textbf{Intégré, déterministe} \\
    \hline
  \end{tabular}
  \vspace{0.2cm}
  \caption{Positionnement comparatif des principales familles de solutions de chargement de données pour l'entraînement à grande échelle.}
  \label{tab:etat_art}
\end{table}

Ce positionnement met en évidence qu'aucune des familles existantes ne combine simultanément les quatre propriétés recherchées : un format binaire à accès direct par projection mémoire, une collation dynamique évitant le gaspillage du padding statique, une coordination déterministe du parallélisme de données, et une absence totale de copie mémoire superflue jusqu'au transfert vers le GPU. Les solutions orientées séquentielles (archives lues en flux continu) sacrifient l'accès aléatoire, indispensable pour certaines stratégies de mélange ou de reprise fine après un point de contrôle. Les solutions orientées vision pure ne traitent pas nativement l'appariement image-texte propre aux VLM. Les solutions en couches, conçues pour la gouvernance et le versionnement de grands corpus, introduisent une couche d'indirection réseau qui pénalise le débit brut. Le chargeur générique de référence, enfin, laisse à la charge de l'utilisateur l'intégralité de la logique de sérialisation efficace, de sorte que son coût réel dépend entièrement de la qualité du code applicatif qui l'entoure code qui, dans la pratique observée, recopie fréquemment les données à chaque étage de transformation.

Cet espace de conception vacant est formulé, pour le reste de ce mémoire, sous la forme de trois questions de recherche directrices :

\begin{itemize}
  \item \textbf{RQ1} : Un mécanisme de lecture zéro-copie par projection mémoire surpasse-t-il, en débit soutenu, une lecture conventionnelle par appels système successifs ?
  \item \textbf{RQ2} : Une stratégie de padding dynamique réduit-elle, conformément à la dépendance quadratique établie au chapitre précédent, le volume de calcul gaspillé par rapport à un padding statique ?
  \item \textbf{RQ3} : La scalabilité du parallélisme de préparation des données suit-elle, en pratique, la loi d'Amdahl théorique, et si oui avec quelle fraction séquentielle effective ?
  \end{itemize}

Ces trois questions structurent directement le protocole d'évaluation présenté en fin de chapitre, et trouveront leur réponse empirique au chapitre consacré aux résultats.

\section{Architecture globale en trois étages}
\label{sec:archi_globale}

L'architecture retenue répond aux exigences dégagées ci-dessus par une séparation stricte en trois étages fonctionnels, illustrée à la figure \ref{fig:archi_trois_etages}.

\begin{itemize}
  \item \textbf{Étage 1 : Prétraitement hors-ligne.} Exécuté une seule fois par corpus, cet étage transforme le jeu de données brut (images et annotations textuelles dans leur format source) en un ensemble de fichiers binaires auto-portants et en un manifeste de coordination. Il est implémenté en Python et exploite le parallélisme multiprocessus décrit au chapitre précédent, dans la mesure où le décodage d'image et la tokenisation les deux opérations les plus coûteuses du pipeline naïf y sont concentrées.
  \item \textbf{Étage 2 : Runtime de chargement.} Exécuté à chaque itération d'entraînement, cet étage ne fait que de la \emph{lecture pure} : aucun décodage d'image compressée, aucune tokenisation à chaud. Il combine la projection mémoire zéro-copie, un thread de préchargement asynchrone, et la collation dynamique en batch, formant le cœur bas niveau implémenté en C++.
  \item \textbf{Étage 3 : Intégration distribuée.} Ce dernier étage assure la coordination entre rangs d'entraînement (partitionnement déterministe des fragments de données) et le transfert des batches assemblés vers la mémoire du GPU par un chemin sans copie intermédiaire côté hôte.
\end{itemize}

Cette séparation offline/online repose sur un principe d'amortissement du coût : le prétraitement, coûteux mais exécuté une seule fois par corpus, paie par avance le coût du décodage et de la tokenisation ; le runtime, exécuté potentiellement des dizaines de fois par échantillon au fil des époques d'entraînement, ne supporte plus que le coût bien moindre de la lecture d'un format déjà normalisé. Un pipeline naïf qui redécoderait chaque image JPEG et retokeniserait chaque texte à chaque époque paierait ce coût élevé de façon répétée, alors que la présente architecture ne le paie qu'une fois.

\begin{figure}[htbp]
  \centering
  \begin{tikzpicture}[
    scale=0.9, transform shape,
    node distance=2cm,
    stage/.style={draw, rectangle, rounded corners, minimum width=5.2cm, minimum height=1.2cm, align=center, font=\small},
    sub/.style={draw, rectangle, minimum width=3.8cm, minimum height=0.7cm, align=center, font=\footnotesize},
    arr/.style={->, thick}
  ]
    \node[stage, fill=gray!10] (s1) {\textbf{Étage 1 — Hors-ligne}\\Décodage, tokenisation,\\écriture binaire (Python)};
    \node[sub, below=0.3cm of s1] (manifest) {Fichiers binaires\\+ manifeste};

    \node[stage, right=1.2cm of s1, fill=gray!10] (s2) {\textbf{Étage 2 — Runtime}\\Projection mémoire, prefetch,\\collation dynamique (C++)};
    \node[sub, below=0.3cm of s2] (queue) {File de prefetch bornée\\($\mu_{\text{loader}} > \lambda_{\text{GPU}}$)};

    \node[stage, right=1.2cm of s2, fill=gray!10] (s3) {\textbf{Étage 3 — Distribué}\\Partitionnement par rang,\\transfert GPU zéro-copie};

    \draw[arr] (s1) -- (s2) node[midway, above] {\footnotesize lecture};
    \draw[arr] (s2) -- (s3) node[midway, above] {\footnotesize batch prêt};
  \end{tikzpicture}
  \vspace{0.5cm}
  \caption{Architecture en trois étages : prétraitement hors-ligne, runtime de chargement, intégration distribuée. La frontière entre l'étage 1 et les étages 2--3 sépare le coût amorti (payé une fois) du coût récurrent (payé à chaque itération).}
  \label{fig:archi_trois_etages}
\end{figure}

La communication entre les étages 1 et 2 se limite à des artefacts sur disque (fichiers binaires et manifeste), sans dépendance d'exécution commune : l'étage 1 peut être exécuté sur une machine distincte de celle qui exécute l'entraînement, ce qui découple les ressources de calcul CPU nécessaires à la préparation des données de celles nécessaires à l'entraînement lui-même.

\subsection{Modèle de stabilité en file d'attente}

Le couple formé par le runtime de chargement (étage 2) et le GPU consommateur peut être modélisé, comme introduit au chapitre précédent, comme un système à un serveur au sens de la théorie des files d'attente : le chargeur constitue le serveur, de taux de service $\mu_{\text{loader}}$ (batches produits par seconde), et le GPU constitue le flux d'arrivée de demandes, de taux $\lambda_{\text{GPU}}$ (batches consommés par seconde). La condition de \textbf{stabilité} du système l'absence de croissance non bornée de la file d'attente, donc l'absence de famine chronique du GPU s'écrit :

\begin{equation}
\label{eq:stabilite}
\mu_{\text{loader}} > \lambda_{\text{GPU}}
\end{equation}

Le facteur d'utilisation $\rho = \lambda_{\text{GPU}} / \mu_{\text{loader}}$ caractérise la marge de sécurité du système : un $\rho$ proche de 1 signifie que le chargeur fonctionne près de sa capacité maximale, rendant le système vulnérable à toute variance transitoire de latence (un seul pic de latence suffisant alors à provoquer une attente du GPU), tandis qu'un $\rho$ significativement inférieur à 1 offre une marge d'absorption du \textit{jitter} sans recourir à une profondeur de prefetch excessive. Ce facteur d'utilisation constitue, avec le débit soutenu lui-même, l'une des grandeurs centrales du protocole d'évaluation défini en fin de chapitre.

\section{Conception du format binaire de sérialisation}
\label{sec:format_binaire}

Le format de sérialisation constitue le point de rencontre entre l'étage hors-ligne, qui l'écrit une seule fois, et le runtime, qui doit pouvoir y accéder de façon aléatoire, en lecture seule, avec un coût d'ouverture minimal. Ce problème est modélisé comme un problème classique de sérialisation à accès aléatoire, dont la solution retenue s'articule en quatre blocs successifs sur le support de stockage : un en-tête de taille fixe, une suite d'enregistrements d'échantillons de taille variable et alignée, une table d'offsets, et un pied de fichier.

\begin{table}[htbp]
  \centering
  \renewcommand{\arraystretch}{1.5}
  \begin{tabular}{|p{3.5cm}|p{3.5cm}|p{8cm}|}
    \hline
    \textbf{Bloc} & \textbf{Taille} & \textbf{Rôle} \\
    \hline
    En-tête & Fixe (quelques dizaines d'octets) & Nombre magique de détection de format, numéro de version, nombre d'échantillons, position de la table d'offsets, dimensions de référence, indicateurs de fonctionnalités (drapeaux) \\
    \hline
    Enregistrements d'échantillons & Variable, alignée & Données de chaque échantillon (image, séquences de tokens, métadonnées), chacune démarrant sur une frontière d'alignement fixe \\
    \hline
    Table d'offsets & Proportionnelle au nombre d'échantillons & Un couple (position, longueur) par échantillon, permettant l'accès direct sans balayage \\
    \hline
    Pied de fichier & Fixe (quelques octets) & Somme de contrôle d'intégrité et marqueur de fin \\
    \hline
  \end{tabular}
  \vspace{0.2cm}
  \caption{Structure conceptuelle du format binaire de sérialisation à accès aléatoire.}
  \label{tab:format_shard}
\end{table}

Chaque champ de l'en-tête répond à un besoin fonctionnel précis. Le nombre magique permet à tout lecteur de rejeter immédiatement un fichier corrompu ou d'un format étranger, sans tenter d'en interpréter le contenu. Le numéro de version permet une évolution ascendante du format : un lecteur récent peut interpréter correctement un fichier ancien, tandis qu'un lecteur ancien rejette explicitement un fichier trop récent plutôt que de l'interpréter de façon erronée. Le champ de drapeaux (\textit{flags}) constitue le point d'extensibilité du format stockage en entiers non signés plutôt qu'en flottant, présence de dimensions par échantillon, compression optionnelle sans nécessiter de nouvelle révision de version pour chaque fonctionnalité additionnelle : un lecteur ignorant un drapeau qu'il ne reconnaît pas peut, selon la sémantique retenue, soit le rejeter explicitement, soit l'ignorer sans danger si la fonctionnalité correspondante est rétrocompatible.

Un choix de conception déterminant concerne la position de la table d'offsets : celle-ci est placée en \emph{fin} de fichier plutôt qu'en tête, immédiatement après le dernier enregistrement d'échantillon. Ce choix permet à l'étage d'écriture de fonctionner en flux continu, sans connaître à l'avance le nombre total d'échantillons que contiendra le fichier : chaque échantillon est écrit dès qu'il est disponible, sa position et sa longueur étant simplement accumulées en mémoire, et ce n'est qu'à la fermeture du fichier que la table complète est sérialisée, son emplacement étant alors reporté dans l'en-tête. Une conception alternative, plaçant la table en tête, aurait exigé soit de connaître $N$ à l'avance, soit de réserver un espace fixe potentiellement insuffisant, soit de réécrire l'en-tête après coup une opération incompatible avec une écriture strictement séquentielle sur certains supports de stockage.

\subsection{Coût d'accès et surcharge d'alignement}

La présence de la table d'offsets transforme la localisation d'un échantillon arbitraire en une opération à coût constant :

\begin{equation}
\label{eq:acces_o1}
T_{\text{accès}}(\text{table}) = O(1) \qquad \text{contre} \qquad T_{\text{accès}}(\text{balayage}) = O(N)
\end{equation}

un lecteur cherchant l'échantillon d'indice $k$ calcule directement l'adresse de son entrée dans la table ($\text{position\_table} + 16k$ pour une entrée de deux champs de 8 octets), lit l'offset et la longueur qui y sont stockés, puis accède directement à l'enregistrement correspondant sans jamais parcourir les enregistrements précédents. Cette propriété est la condition nécessaire pour qu'un mélange aléatoire de l'ordre de lecture (accès à des indices non contigus) ne dégrade pas le débit du lecteur par rapport à une lecture strictement séquentielle, hormis l'effet, déjà quantifié au chapitre précédent, de la perte de localité spatiale sur la politique de lecture anticipée du système d'exploitation.

L'alignement de chaque enregistrement sur une frontière fixe motivé, comme établi précédemment, par la largeur des lignes de cache du processeur introduit un espace de rembourrage (\textit{padding}) résiduel entre la fin utile d'un enregistrement et la frontière d'alignement suivante. En supposant que cette longueur de rembourrage suit une distribution uniforme sur l'intervalle $[0, \text{align}-1]$ hypothèse raisonnable dès lors que la taille des enregistrements n'est pas elle-même corrélée à l'alignement retenu la surcharge moyenne attendue par échantillon s'écrit :

\begin{equation}
\label{eq:surcharge_alignement}
\mathbb{E}[\text{surcharge}] = \frac{\text{align}}{2}
\end{equation}

soit, pour un alignement de 64 octets, une surcharge moyenne de 32 octets par échantillon négligeable devant la taille typique d'un enregistrement (plusieurs kilooctets pour une image), mais non nulle, et proportionnelle au nombre total d'échantillons du corpus. La table d'offsets elle-même contribue à ce coût fixe de stockage à un taux de 16 octets par échantillon, indépendamment de la taille des données utiles un surcoût structurel acceptable au regard du gain d'accès à coût constant qu'il permet.

\section{Modélisation du pipeline de préchargement asynchrone}
\label{sec:prefetch}

Le runtime de chargement (étage 2) est spécifié comme un automate producteur-consommateur à deux threads. Un premier flux d'exécution, le \textbf{producteur}, projette en mémoire les fichiers binaires successifs, en extrait les enregistrements d'échantillons via la table d'offsets, et les pousse dans une structure de communication sans verrou introduite au chapitre précédent. Un second flux d'exécution, le \textbf{consommateur}, retire de cette structure les batches déjà assemblés et les met à disposition de la boucle d'entraînement.

La profondeur de la file de communication entre ces deux flux le \textbf{prefetch} est le paramètre central de dimensionnement de cet étage, dont les deux régimes de mauvais dimensionnement sont symétriques et également indésirables : une profondeur trop faible expose le système à la famine du GPU décrite au chapitre précédent dès qu'un seul batch tarde à être produit, tandis qu'une profondeur excessive immobilise, sans bénéfice supplémentaire, une mémoire proportionnelle au nombre de batches maintenus en réserve. Sous l'hypothèse d'une file bornée de capacité $K$, le débit stable du système correspond au minimum des taux de service des deux extrémités :

\begin{equation}
\label{eq:debit_stable}
\mu_{\text{stable}} = \min(\mu_{\text{prod}}, \mu_{\text{cons}})
\end{equation}

et, dans le régime déséquilibré où l'un des deux taux domine nettement l'autre, la file tend à se remplir (ou à se vider) presque intégralement, de sorte que le temps de réponse moyen d'un élément dans le système s'approche de $K/2$ une approximation directement dérivée de la loi de Little $L = \lambda W$ introduite précédemment, appliquée ici au régime où la file oscille entre un état presque plein et un état presque vide plutôt que de se stabiliser à une occupation intermédiaire constante.

\subsection{Localité de lecture et mélange par fenêtre glissante}

Pour préserver la localité spatiale exploitée par la projection mémoire, le producteur ne traite qu'un seul fichier binaire à la fois, dans un ordre séquentiel, plutôt que d'entrelacer les accès entre plusieurs fichiers simultanément un choix qui maximise l'efficacité de la lecture anticipée du système d'exploitation au prix d'une contrainte sur la stratégie de mélange des données.

Un mélange strictement uniforme sur l'ensemble du corpus exigerait de matérialiser en mémoire un index global permutable, ce qui contredirait la propriété de mémoire bornée établie au chapitre précédent. La stratégie retenue est un compromis classique en présence de données trop volumineuses pour tenir en mémoire : un \textbf{tampon de mélange} (\textit{shuffle buffer}) de taille bornée et configurable reçoit les échantillons au fil de leur lecture séquentielle, et restitue, à chaque demande du consommateur, un échantillon tiré aléatoirement parmi ceux actuellement présents dans le tampon, remplacé aussitôt par le prochain échantillon lu séquentiellement. Le générateur pseudo-aléatoire sous-jacent est initialisé de façon déterministe à partir d'une graine combinée à l'indice de l'époque courante, garantissant qu'une même exécution reproduit exactement le même ordre de présentation des données d'une exécution à l'autre, tout en produisant un ordre différent à chaque époque.

Ce mécanisme n'approxime un mélange véritablement uniforme que dans la limite où la taille du tampon devient grande devant la structure locale du corpus (par exemple, le regroupement d'échantillons proches dans le fichier source) : plus le tampon est profond, plus deux échantillons consécutifs dans le fichier source ont de chances d'être physiquement séparés dans l'ordre de présentation résultant, au prix d'une empreinte mémoire proportionnelle à la taille du tampon. Un tampon de taille nulle dégénère en une lecture purement séquentielle, sans aucun mélange un régime parfois recherché délibérément pour le déboguage ou la reproduction exacte d'un ordre de traitement.

\section{Stratégie de padding dynamique et normalisation différée}
\label{sec:padding_dynamique}

La décision de conception la plus directement dérivée du cadre théorique du chapitre précédent consiste à ne jamais matérialiser de padding au moment de l'écriture des données : chaque échantillon conserve, jusqu'à l'instant de la collation en batch, ses dimensions réelles propres.

Concrètement, chaque enregistrement d'échantillon porte un court préfixe encodant ses dimensions effectives (hauteur et largeur réelles de l'image après redimensionnement, en plus des dimensions d'origine avant redimensionnement, conservées à des fins de traçabilité). Au moment de la collation, pour un batch de $B$ échantillons, le runtime détermine les dimensions maximales effectivement présentes,

\begin{equation}
\label{eq:dims_batch}
H_{\text{batch}} = \max_i h_i, \qquad W_{\text{batch}} = \max_i w_i, \qquad L_{\text{batch}} = \max_i \ell_i
\end{equation}

où $\ell_i$ désigne la longueur de la séquence de tokens de l'échantillon $i$, puis assemble le batch en un unique passage sur les $B$ échantillons — une opération en $O(B)$, indépendante du nombre total d'échantillons du corpus. La normalisation de type ImageNet, $(x/255 - \mu)/\sigma$, est appliquée à ce moment précis, sur des valeurs initialement stockées en entiers non signés puis converties en flottant : cette opération, un produit-somme élémentaire appliqué indépendamment à chaque pixel, est directement vectorisable par les unités SIMD, la commutation entre stockage compact et calcul flottant n'intervenant qu'au tout dernier instant du pipeline.

\subsection{Gain mémoire et neutralité de la normalisation différée}

Le gain mémoire apporté par cette stratégie, relativement à une politique de padding statique qui matérialiserait systématiquement chaque image à une dimension fixe $H_{\text{fixe}} \times W_{\text{fixe}}$, s'exprime pour un batch de $B$ échantillons à $C$ canaux comme :

\begin{equation}
\label{eq:gain_memoire}
G_{\text{mémoire}} = 1 - \frac{\sum_{i=1}^{B} C \cdot h_i \cdot w_i}{B \cdot C \cdot H_{\text{fixe}} \cdot W_{\text{fixe}}}
\end{equation}

un gain qui croît avec l'hétérogénéité des tailles présentes dans le corpus et se rapproche de zéro lorsque toutes les images du corpus partagent déjà des dimensions proches de $H_{\text{fixe}} \times W_{\text{fixe}}$ cas dans lequel le padding dynamique n'apporte, par construction, qu'un bénéfice marginal.

La validité de ce report de la normalisation à l'étape de collation repose sur un argument de commutativité algébrique déjà esquissé au chapitre précédent, qu'il convient ici de formaliser pleinement dans le contexte de l'assemblage d'un batch. La normalisation étant une transformation affine appliquée indépendamment à chaque pixel, elle commute avec l'opération de concaténation qui assemble les échantillons individuels en un tenseur de batch : normaliser chaque image séparément puis les concaténer produit un résultat rigoureusement identique à celui obtenu en concaténant d'abord les images brutes puis en appliquant la normalisation au tenseur résultant, la concaténation ne mélangeant les valeurs d'aucun pixel entre elles. Cette propriété de linéarité est précisément ce qui autorise à différer la normalisation du moment de l'écriture (un pixel à la fois, potentiellement des millions de fois au fil des époques si elle était répétée) au moment de la collation (également un pixel à la fois, mais une seule fois par lecture effective, et vectorisable sur l'ensemble du batch plutôt que pixel par pixel).

\section{Méthodologie d'évaluation et protocole expérimental}
\label{sec:methodologie}

Conformément aux bonnes pratiques de conception expérimentale, les métriques, les bases de comparaison et les menaces à la validité des mesures sont définies préalablement à toute campagne de mesure, avant même l'implémentation des instruments de mesure eux-mêmes une discipline destinée à éviter le biais consistant à ajuster a posteriori les métriques rapportées en fonction des résultats obtenus.

\begin{table}[htbp]
  \centering
  \renewcommand{\arraystretch}{1.5}
  \begin{tabular}{|p{4.5cm}|p{5.5cm}|p{4.5cm}|}
    \hline
    \textbf{Métrique} & \textbf{Définition} & \textbf{Question de recherche associée} \\
    \hline
    Débit soutenu & Échantillons et batches traités par seconde en régime stationnaire & RQ1 \\
    \hline
    Latence par batch & Moyenne, médiane (P50) et 99\textsuperscript{e} centile (P99) du temps de production d'un batch & RQ1 \\
    \hline
    Pic de mémoire résidente & Mémoire maximale occupée durant l'exécution & RQ3 \\
    \hline
    Empreinte disque & Taille du corpus sérialisé rapportée à sa taille source & RQ2 \\
    \hline
    Proxy de calcul gaspillé & FLOPs d'attention estimés sous padding statique vs dynamique & RQ2 \\
    \hline
    Accélération et efficacité parallèle & $S(n)$ et $E(n) = S(n)/n$ en fonction du nombre de workers & RQ3 \\
    \hline
  \end{tabular}
  \vspace{0.2cm}
  \caption{Métriques retenues pour le protocole d'évaluation empirique.}
  \label{tab:metriques}
\end{table}

Le protocole retenu impose l'usage exclusif de corpus réels le rejet explicite de données synthétiques comme base d'évaluation principale afin que la distribution mesurée des tailles d'image et des longueurs de texte reflète les conditions réelles d'un entraînement, et non un cas favorable arbitrairement construit. Une graine pseudo-aléatoire fixe est utilisée systématiquement pour toute opération impliquant un tirage (échantillonnage, mélange), garantissant la reproductibilité intégrale des mesures d'une exécution à l'autre. Une phase de chauffe (\textit{warmup}), dont la durée est exclue du chronométrage, précède chaque campagne de mesure afin de neutraliser les effets transitoires de premier accès remplissage initial des caches mémoire et disque, compilation différée éventuelle qui biaiseraient sinon les toutes premières mesures vers le haut. Chaque configuration est répétée sur un nombre suffisant de batches pour permettre le calcul d'une statistique robuste (médiane plutôt que moyenne arithmétique, moins sensible aux valeurs aberrantes) ainsi que des centiles de latence, et l'environnement matériel (processeur, mémoire vive, type de support de stockage, système d'exploitation) est systématiquement documenté aux côtés des résultats, ces derniers n'étant significatifs que rapportés à la configuration matérielle qui les a produits.

Pour une estimation du débit moyen fondée sur $n$ mesures indépendantes de latence par batch, de moyenne empirique $\bar{x}$ et d'écart-type empirique $s$, l'intervalle de confiance à 95\% est approximé, sous hypothèse de normalité asymptotique du grand échantillon, par :

\begin{equation}
\label{eq:ic_95}
\bar{x} \pm 1{,}96 \, \frac{s}{\sqrt{n}}
\end{equation}

un intervalle qui se resserre en $O(1/\sqrt{n})$ à mesure que le nombre de répétitions augmente, justifiant le choix d'un nombre de batches mesurés suffisamment élevé pour que la largeur de cet intervalle reste négligeable devant l'écart attendu entre les configurations comparées.

\bigskip
\noindent Ce chapitre a traduit le cadre théorique du chapitre précédent en une architecture concrète à trois étages, un format de sérialisation à accès aléatoire, un pipeline de préchargement asynchrone, une stratégie de padding dynamique, et un protocole d'évaluation défini par avance autour de trois questions de recherche. Le chapitre suivant décrira l'implémentation logicielle qui matérialise ces choix de conception.

% PAGE DE CHAPITRE STYLISÉE AVEC LOGOS ET TITRE CENTRÉ
\begin{titlepage}
  \thispagestyle{empty} % Supprime les en-têtes/pieds pour cette page

  % LOGOS HAUT DE PAGE
  \begin{minipage}{0.48\textwidth}
  \includegraphics[width=7cm]{Images_/logo_ensam_mek.png}\\[0.2cm]
  \textbf{}
\end{minipage}
\hfill
\begin{minipage}{0.48\textwidth}
  \raggedleft
  % Pense à ajouter le logo de AI Movement dans ton dossier Images_
  \includegraphics[width=5.5cm]{Images_/aimovement_logo.png}\\[0.4cm]
  \textbf{} \\
\end{minipage}

  % ESPACE VIDE POUR PLACEMENT VERTICAL AU CENTRE
  \vspace*{5cm}

  % TITRE AU CENTRE
  \begin{center}
    \rule{\linewidth}{0.5pt}\vspace{0.7cm}

    {\Huge\bfseries\color{blue}
    Chapitre 4 :\\[0.5em]
    Implémentation : Prétraitement Python et Génération des Shards}

    \vspace{0.7cm}
    \rule{\linewidth}{0.5pt}
  \end{center}

  \vfill

\end{titlepage}

\chapter{Implémentation : Prétraitement Python et Génération des Shards}
\label{chap:implementation_python}

\section{Système de configuration déclaratif et schéma de données}
\label{sec:config_schema}

Le package de prétraitement repose sur une séparation stricte entre la configuration et le code. Aucun paramètre de traitement n'est codé en dur dans la logique applicative. Tout comportement observable du pipeline (dimensions cibles, stratégie de redimensionnement, modèle de tokeniseur, taille des shards) se règle exclusivement par l'édition d'un fichier de configuration au format YAML.

Cette configuration est chargée dans une hiérarchie de classes de données Python immuables (des \textit{dataclasses} figées), organisée en quatre sous configurations spécialisées, elles mêmes assemblées dans une configuration de niveau supérieur. La classe \texttt{DatasetConfig} décrit la source des données et le mappage des colonnes. La classe \texttt{ImageConfig} décrit les paramètres de prétraitement d'image. La classe \texttt{TokenizerConfig} décrit les paramètres de tokenisation. La classe \texttt{ShardConfig} décrit les paramètres d'écriture des fichiers binaires. La classe \texttt{PipelineConfig} rassemble ces quatre sous configurations et ajoute les paramètres globaux de l'exécution (nombre de workers, stratégie de mélange, niveau de journalisation, graine aléatoire).

L'immutabilité de ces classes n'est pas un choix esthétique. Dans un pipeline multiprocessus, une configuration partagée entre plusieurs processus enfants doit être sérialisée puis reconstruite dans chaque processus. Une configuration mutable exposerait le risque qu'un worker modifie accidentellement son état local sans que cette modification se propage aux autres, produisant un comportement incohérent d'un worker à l'autre. Le figement de la structure élimine cette classe d'erreurs par construction.

Le chargement d'un fichier de configuration s'accompagne d'une validation systématique. L'absence du fichier, une syntaxe YAML invalide, ou une valeur hors de l'ensemble autorisé (par exemple une stratégie de mélange non reconnue) déclenchent une erreur explicite au moment du chargement plutôt qu'une défaillance différée et difficile à diagnostiquer en cours d'exécution.

\begin{table}[htbp]
  \centering
  \renewcommand{\arraystretch}{1.5}
  \begin{tabular}{|p{2.5cm}|p{5cm}|p{6.5cm}|}
    \hline
    \textbf{Section} & \textbf{Rôle} & \textbf{Paramètres principaux} \\
    \hline
    \texttt{dataset} & Source des données et mappage & Identifiant du jeu, découpage (\textit{split}), mappage canonique, limite d'échantillons \\
    \hline
    \texttt{image} & Prétraitement des images & Dimensions cibles, stratégie de redimensionnement, moyennes et écarts types \\
    \hline
    \texttt{tokenizer} & Tokenisation du texte & Modèle, longueur maximale, stratégie de rembourrage, troncature \\
    \hline
    \texttt{shard} & Écriture des fichiers binaires & Répertoire, taille cible, nombre d'échantillons max, alignement, version \\
    \hline
    \textit{Niveau global} & Comportement du pipeline & Nombre de workers, stratégie de mélange, graine, niveau de journalisation \\
    \hline
  \end{tabular}
  \vspace{0.2cm}
  \caption{Les cinq sections de la configuration déclarative et leurs paramètres principaux.}
  \label{tab:config_sections}
\end{table}

Un mécanisme de mappage déclaratif des colonnes rend le pipeline indépendant de la structure interne du jeu de données source. Plutôt que d'écrire du code spécifique à chaque jeu de données pour extraire l'image, la question et la réponse, la configuration associe des noms canoniques fixes (\texttt{image}, \texttt{question}, \texttt{answer}, \texttt{id}, \texttt{reasoning}) aux noms de colonnes réels du jeu de données choisi. Changer de source de données revient alors à éditer quelques lignes de configuration, sans toucher au code d'ingestion.

En sortie de l'étape d'ingestion, chaque enregistrement du jeu de données source est converti en une instance d'un schéma canonique unique, la classe \texttt{VQASample}. Cette classe porte sept champs : un identifiant unique, les octets bruts de l'image dans son encodage d'origine (non décodée à ce stade, le décodage étant différé à l'étape de normalisation), la largeur et la hauteur d'origine de l'image, le texte de la question, le texte de la réponse, et un dictionnaire libre de métadonnées auxiliaires destiné aux champs non essentiels au schéma (par exemple un raisonnement intermédiaire fourni par certains jeux de données).

Cette classe expose une méthode \texttt{validate()}, constituant la première ligne de défense contre les données corrompues avant que celles ci n'entrent dans le reste du pipeline. Elle vérifie que l'identifiant n'est pas vide, que les octets d'image sont non vides et de type correct, que les dimensions annoncées sont strictement positives, et que la question comme la réponse contiennent un texte non vide après élimination des espaces superflus. Une méthode distincte et plus coûteuse, \texttt{validate\_image\_decodable()}, tente le décodage effectif des octets d'image par la bibliothèque graphique sous jacente, une vérification volontairement réservée aux contrôles de qualité en amont plutôt qu'appliquée systématiquement à chaque accès, en raison de son coût.

\section{Traitement d'image : redimensionnement préservant le ratio d'aspect}
\label{sec:image_processing}

Le module de traitement d'image expose en réalité deux chemins distincts, correspondant aux deux versions du format binaire introduites au chapitre précédent. Le premier chemin, employé par le format v1, produit une image normalisée de dimensions fixes selon l'une de trois stratégies configurables. Le second chemin, employé par le format v2, préserve systématiquement le ratio d'aspect de l'image et ne matérialise jamais de padding, conformément à la stratégie de padding dynamique développée au chapitre précédent.

\subsection{Les trois stratégies du chemin v1}

Trois stratégies de redimensionnement sont disponibles pour le chemin à dimensions fixes. Le \textbf{redimensionnement avec remplissage} ramène le côté le plus long de l'image à la dimension cible, redimensionne l'autre côté proportionnellement, puis centre le résultat sur une toile de la taille cible en comblant l'espace restant avec une valeur de remplissage constante. Le \textbf{recadrage centré} procède à l'inverse, en ramenant le côté le plus court à la dimension cible puis en découpant l'excédent du côté le plus long, centré. Le \textbf{redimensionnement direct} force les deux dimensions à leur valeur cible sans considération du ratio d'aspect d'origine.

\begin{table}[htbp]
  \centering
  \renewcommand{\arraystretch}{1.5}
  \begin{tabular}{|p{4cm}|p{2.8cm}|p{4.2cm}|p{4.2cm}|}
    \hline
    \textbf{Stratégie} & \textbf{Ratio d'aspect} & \textbf{Perte d'information} & \textbf{Effet visuel} \\
    \hline
    Redimensionnement avec remplissage & Préservé & Aucune (bandes de remplissage ajoutées) & Bandes uniformes visibles sur un des deux bords \\
    \hline
    Recadrage centré & Préservé localement & Coupe des bords, biaisée vers le centre & Aucune bande, mais perte du contenu périphérique \\
    \hline
    Redimensionnement direct & Non préservé & Aucune perte de pixels, mais distorsion géométrique & Image étirée ou compressée selon un axe \\
    \hline
  \end{tabular}
  \vspace{0.2cm}
  \caption{Comparaison des trois stratégies de redimensionnement du chemin à dimensions fixes.}
  \label{tab:resize_strategies}
\end{table}

Le redimensionnement avec remplissage est retenu comme stratégie par défaut du chemin v1. Le recadrage centré est écarté comme choix par défaut car il élimine potentiellement une partie du contenu informatif de l'image, en particulier lorsque l'élément pertinent pour répondre à la question associée ne se trouve pas au centre du cadre, un risque significatif dans un contexte de question réponse visuelle où la zone d'intérêt peut se situer n'importe où dans l'image. Le redimensionnement direct est écarté car il introduit une distorsion géométrique susceptible de dégrader la reconnaissance de formes par l'encodeur visuel, en particulier pour des ratios d'aspect très éloignés du carré.

\subsection{Le chemin v2 à ratio d'aspect préservé}

Le chemin v2 n'admet qu'une seule politique. Le côté le plus long de l'image est ramené à une dimension maximale configurable, et l'autre côté est calculé pour préserver exactement le ratio d'aspect d'origine, sans jamais matérialiser de bande de remplissage. Le facteur d'échelle appliqué s'écrit :

\begin{equation}
\label{eq:facteur_echelle}
s = \min\left(1,\ \frac{D}{\max(h, w)}\right)
\end{equation}

où $D$ est la dimension maximale configurée et $h$, $w$ les dimensions d'origine de l'image. La présence du terme $\min(1, \cdot)$ traduit une décision d'ingénierie importante. Une image déjà plus petite que la dimension maximale n'est jamais agrandie. Seul un sur échantillonnage vers le bas est appliqué, jamais un sur échantillonnage vers le haut, ce dernier n'apportant aucune information supplémentaire et gaspillant du calcul et de l'espace de stockage pour un gain nul. Les nouvelles dimensions résultent de $h' = h \cdot s$ et $w' = w \cdot s$, arrondies à l'entier inférieur.

L'interpolation employée par défaut est bicubique. Ce choix repose sur un compromis entre qualité et coût de calcul : l'interpolation bicubique, dont le noyau de convolution s'étend sur un voisinage de seize pixels contre quatre pour l'interpolation bilinéaire, produit un résultat visuellement plus lisse et réduit les artefacts d'aliasing lors d'une réduction significative de taille, au prix d'un coût de calcul supérieur. Les interpolations bilinéaire, de Lanczos et par plus proche voisin restent disponibles par configuration pour des cas d'usage où le coût de calcul prime sur la qualité de rééchantillonnage, la dernière étant réservée à des cas particuliers où toute interpolation doit être évitée (données catégorielles encodées en pixels, par exemple).

La conversion vers l'espace colorimétrique cible (RVB par défaut) précède le redimensionnement. L'image résultante est ensuite transposée du format hauteur largeur canaux vers le format canaux hauteur largeur, attendu par les encodeurs visuels de type Transformer, et conservée en entiers non signés sur huit bits, sans normalisation. La normalisation est, comme établi au chapitre précédent, différée jusqu'à l'étape de collation du batch.

\section{Tokenisation et alignement des séquences textuelles}
\label{sec:tokenisation}

La tokenisation du texte s'appuie sur un encapsulage léger autour des tokeniseurs de l'écosystème HuggingFace, dont le comportement est entièrement piloté par la configuration (identifiant du modèle de tokeniseur, longueur maximale, stratégie de rembourrage, activation de la troncature). Le chargement du modèle de tokeniseur, une opération de l'ordre de la seconde, n'est effectué qu'une seule fois par instance et réutilisé pour tous les appels suivants.

Une attention particulière est portée à la robustesse face à des modèles de tokeniseur dont la configuration native ne définit pas de jeton de remplissage. Dans ce cas, le jeton de fin de séquence est utilisé comme substitut, avec journalisation explicite de ce choix, plutôt que de laisser l'opération de rembourrage échouer silencieusement ou lever une exception peu explicite plus loin dans le pipeline.

Chaque échantillon produit quatre tableaux distincts, obtenus par deux appels indépendants à la fonction de tokenisation : un pour la question, un pour la réponse. Chaque appel retourne un tableau d'identifiants de tokens et un masque d'attention de même longueur, ce dernier valant un pour chaque position occupée par un jeton réel et zéro pour chaque position de remplissage. Avec la stratégie de rembourrage par défaut (remplissage jusqu'à la longueur maximale configurée), les quatre tableaux ont systématiquement la même longueur fixe au moment de l'écriture, quelle que soit la longueur réelle du texte source.

Le masque d'attention porte ainsi une information qui n'est pas redondante avec la longueur fixe du tableau qui le contient : il encode la longueur réelle de la séquence utile, récupérable en sommant ses valeurs ou en repérant la position du dernier jeton actif. Cette information est indispensable au padding dynamique par batch mis en œuvre au moment de la lecture : même si chaque échantillon est stocké à une longueur fixe, le runtime peut, en lisant le masque, déterminer la longueur effective de chaque séquence du batch courant et n'en conserver que le maximum réellement utile, plutôt que de systématiquement traiter la longueur maximale configurée pour l'ensemble du corpus.

\section{Écriture du format binaire : alignement, table des offsets et CRC32}
\label{sec:ecriture_binaire}

L'écriture des échantillons prétraités dans le format binaire décrit au chapitre précédent est prise en charge par une classe unique, responsable de l'intégralité du cycle de vie d'un fichier de shard. Cette classe expose une interface volontairement réduite à quatre opérations publiques, résumées au tableau \ref{tab:shard_writer_api}.

\begin{table}[htbp]
  \centering
  \renewcommand{\arraystretch}{1.5}
  \begin{tabular}{|p{3.5cm}|p{11.5cm}|}
    \hline
    \textbf{Opération} & \textbf{Comportement} \\
    \hline
    Ouverture & Crée le fichier et écrit immédiatement un en-tête de soixante-quatre octets, avec le compte d'échantillons et la position de la table des offsets provisoirement mis à zéro \\
    \hline
    Ajout d'un échantillon & Aligne la position d'écriture courante sur la frontière configurée, sérialise l'ensemble des champs de l'échantillon en un seul bloc contigu, écrit ce bloc (éventuellement compressé), et enregistre son décalage et sa longueur en mémoire \\
    \hline
    Vérification de rotation & Signale si la taille courante ou le nombre d'échantillons courants a atteint l'un des deux seuils configurés \\
    \hline
    Fermeture & Aligne, écrit la table des offsets accumulée, revient en arrière dans le fichier pour renseigner le compte d'échantillons et la position de la table, calcule la somme de contrôle sur l'intégralité du contenu, puis écrit le pied de fichier \\
    \hline
  \end{tabular}
  \vspace{0.2cm}
  \caption{Cycle de vie de l'écrivain de shards.}
  \label{tab:shard_writer_api}
\end{table}

L'écriture de chaque échantillon commence par un alignement sur la frontière configurée (soixante quatre octets par défaut), motivé par les considérations de cache mémoire établies au chapitre précédent. Sous l'hypothèse d'une longueur de rembourrage uniformément distribuée sur l'intervalle possible, la surcharge moyenne attendue par échantillon vaut la moitié de l'alignement, soit trente deux octets pour un alignement de soixante quatre. L'ensemble des champs de l'échantillon (préfixe de dimensions et données d'image pour le format v2, quatre tableaux de tokens, longueur et contenu des métadonnées sérialisées en JSON) est assemblé en un unique bloc contigu avant écriture, plutôt que par une succession de petites écritures séparées, ce qui limite le nombre d'appels système et permet, lorsque la compression est activée, de compresser l'échantillon dans son intégralité en une seule opération plutôt que champ par champ.

Une option de compression légère peut être activée par configuration. Lorsqu'elle l'est, chaque bloc d'échantillon sérialisé est compressé avant écriture, précédé d'un entier indiquant sa taille non compressée, nécessaire à la décompression côté lecture. Cette option reste désactivée par défaut, le format privilégiant la vitesse de lecture côté runtime C++ plutôt que la compacité du stockage sur disque.

La rotation d'un fichier de shard vers le suivant obéit à une conception à deux seuils indépendants, combinés par une disjonction logique : un fichier se ferme dès que sa taille atteint le seuil configuré en mégaoctets, \emph{ou} dès que son nombre d'échantillons atteint le seuil configuré, le premier des deux critères satisfait déclenchant la rotation. Cette double condition permet de garantir simultanément une taille de fichier prévisible (utile pour la répartition des lectures entre processus) et une limite haute sur le nombre d'échantillons par fichier (utile pour borner le temps de construction d'une table des offsets côté lecture), sans qu'un corpus composé d'échantillons de taille très variable ne fasse dériver l'un des deux critères de façon incontrôlée.

À la fermeture, la table des offsets accumulée en mémoire durant l'écriture est sérialisée à la position courante, chaque entrée occupant seize octets (un décalage et une longueur, chacun sur huit octets). Le fichier est alors rouvert en position pour renseigner rétroactivement, dans l'en tête déjà écrit, le nombre final d'échantillons et la position effective de la table. La somme de contrôle CRC32, calculée sur l'intégralité du contenu précédant le pied de fichier, est alors calculée en une seule passe et écrite en fin de fichier, accompagnée du marqueur de fin.

La vérification de cette somme de contrôle est effectuée symétriquement à la lecture, aussi bien par un lecteur Python léger utilisé pour l'inspection et les tests que par le lecteur C++ du runtime, les deux implémentations calculant le même CRC32 selon la norme ISO 3309, garantissant qu'un fichier corrompu est détecté quel que soit le lecteur qui l'ouvre. Une suite de tests de bout en bout dédiée écrit un ensemble d'échantillons volontairement construits pour couvrir des cas limites (dimensions extrêmes, texte vide après troncature, métadonnées imbriquées), puis relit ce même fichier successivement par le lecteur Python et par le lecteur C++, en exigeant une égalité bit à bit de chaque champ relu avec la valeur d'origine. Cette vérification de bout en bout constitue la garantie ultime que les deux implémentations du format, écrites indépendamment dans deux langages différents, restent parfaitement compatibles.

\section{Pipeline multiprocessus en flux à mémoire bornée}
\label{sec:pipeline_streaming}

Cette section porte sur la partie la plus algorithmique du chapitre : l'orchestration qui transforme un jeu de données ingéré en mémoire en un ensemble de fichiers de shards sur disque, tout en garantissant une empreinte mémoire indépendante de la taille totale du corpus.

\subsection{Découpage en blocs et réutilisation du pool de workers}

Le jeu de données ingéré est découpé en blocs successifs (\textit{chunks}) dont la taille correspond au seuil configuré de nombre d'échantillons par shard, ou à une valeur par défaut de cinq cents échantillons lorsque ce seuil n'est pas fixé. Un unique pool de processus est créé avant le traitement du premier bloc, et réutilisé pour l'ensemble des blocs suivants, plutôt que reconstruit à chaque bloc. Ce choix évite de payer, pour chaque bloc, le coût non négligeable du chargement d'un modèle de tokeniseur dans chaque processus fils.

Le chargement du tokeniseur dans les processus fils, plutôt que dans le processus parent suivi d'un partage, s'explique par une contrainte technique des tokeniseurs de l'écosystème HuggingFace, qui ne garantissent pas un comportement correct après duplication par la primitive \textit{fork} du système d'exploitation. Chaque processus fils du pool initialise donc son propre exemplaire de tokeniseur une seule fois, au démarrage, via une fonction d'initialisation dédiée exécutée automatiquement par le pool.

Pour chaque bloc, les échantillons sont distribués aux processus du pool, chacun exécutant indépendamment l'étape de normalisation d'image et de tokenisation. Toute erreur affectant un échantillon individuel (image indécodable, texte manquant, dépassement d'une contrainte de validation) est capturée localement dans le processus qui la rencontre : l'échantillon fautif est simplement écarté, un compteur d'erreurs est incrémenté, et le traitement du reste du bloc se poursuit sans interruption. Le pipeline complet n'échoue donc jamais à cause d'un unique échantillon corrompu, un comportement essentiel face à des corpus réels dont un faible pourcentage d'entrées est presque toujours défectueux.

\subsection{Mélange local et bornitude de la mémoire}

Lorsque la stratégie de mélange local est activée par configuration, chaque bloc traité est mélangé indépendamment avant écriture, au moyen d'un unique générateur pseudo aléatoire initialisé une fois à partir de la graine globale de configuration puis réutilisé, dans le processus principal, pour chaque bloc successif. Ce mélange s'applique uniquement à l'intérieur d'un même bloc : il ne redistribue jamais d'échantillons entre deux blocs distincts, ce qui préserve la propriété centrale de cette architecture, à savoir qu'à tout instant seul un bloc est intégralement matérialisé en mémoire.

Une fois un bloc écrit sur disque, ses données prétraitées sont explicitement libérées et un cycle de collecte des objets non référencés est déclenché avant le passage au bloc suivant. Cette libération explicite garantit que la mémoire de pointe du pipeline reste bornée par la taille d'un seul bloc, indépendamment du nombre total de blocs à traiter, conformément au modèle établi au chapitre précédent.

\begin{figure}[htbp]
  \centering
  \begin{tikzpicture}[
    node distance=0.9cm,
    box/.style={draw, rectangle, rounded corners, minimum width=3.6cm, minimum height=0.9cm, align=center, font=\small},
    arr/.style={->, thick}
  ]
    \node[box] (chunk) {Bloc borné de $N$ échantillons};
    \node[box, below=of chunk] (pool) {Pool de processus réutilisé\\(tokeniseur initialisé une fois)};
    \node[box, below=of pool] (shuffle) {Mélange local optionnel\\(RNG seedée, une passe)};
    \node[box, below=of shuffle] (write) {Écriture séquentielle\\du fichier de shard};
    \node[box, below=of write] (flush) {Libération explicite\\de la mémoire du bloc};

    \draw[arr] (chunk) -- (pool);
    \draw[arr] (pool) -- (shuffle);
    \draw[arr] (shuffle) -- (write);
    \draw[arr] (write) -- (flush);
    \draw[arr] (flush.east) to[out=0, in=0, looseness=1.6] node[right, align=left, font=\footnotesize] {bloc suivant} (chunk.east);
  \end{tikzpicture}
  \caption{Cycle de traitement d'un bloc dans le pipeline en flux. Le pool de processus est créé une seule fois, hors de cette boucle, et réutilisé à chaque itération.}
  \label{fig:pipeline_flux}
\end{figure}

La garantie de bornitude peut s'exprimer, en première approximation, comme une mémoire de pointe proportionnelle à la taille d'un bloc et au nombre de processus du pool, augmentée d'une composante fixe liée au processus d'écriture lui même. Cette propriété est vérifiée empiriquement par une batterie de six tests dédiés au pipeline en flux, couvrant respectivement la production effective de fichiers de shards, l'exactitude du nombre de fichiers produits, l'égalité du nombre total d'échantillons traités entre le mode en flux et un mode séquentiel de référence, la tolérance aux erreurs sur des échantillons individuels sans interruption du pipeline, le bon fonctionnement du mélange local, et enfin la bornitude effective de la mémoire, ce dernier test exécutant le pipeline sur un corpus de taille $N$ puis sur un corpus de taille $3N$, et vérifiant que le pic de mémoire tracée pour le second run reste dans un rapport raisonnable avec le premier plutôt que de croître proportionnellement à la taille du corpus.

\section{Manifeste distribué et attribution déterministe des shards}
\label{sec:manifeste_python}

Une fois l'ensemble des fichiers de shards produits, un manifeste au format JSON est construit par balayage du répertoire de sortie. Cette étape ouvre chaque fichier, en lit uniquement les soixante quatre premiers octets pour en extraire le nombre d'échantillons et vérifier le nombre magique attendu, mesure sa taille sur disque, et calcule sa somme de contrôle CRC32 en excluant les douze derniers octets correspondant au pied de fichier déjà écrit par l'étape d'écriture. Le manifeste résultant regroupe, pour chaque fichier, son chemin relatif, son nombre d'échantillons, sa taille en octets et sa somme de contrôle hexadécimale, ainsi que des totaux agrégés sur l'ensemble du corpus et un horodatage de construction.

Ce manifeste constitue le contrat de coordination entre l'étape de préparation des données et l'étape d'entraînement distribué. Toute application souhaitant participer à l'entraînement peut le charger, sans avoir besoin de parcourir elle même le système de fichiers ni de rouvrir chaque fichier de shard pour en connaître le contenu.

L'attribution des fichiers de shards à un rang donné suit exactement le schéma en répartition cyclique formalisé au chapitre précédent : le fichier d'indice $i$ dans l'ordre du manifeste est attribué au rang $i \bmod W$, où $W$ désigne le nombre total de rangs. Cette fonction est pure, au sens où elle ne dépend d'aucun état mutable ni d'aucune communication entre processus : chaque rang peut recalculer localement, à partir du seul manifeste partagé, la liste des fichiers qui lui reviennent, sans coordination réseau. Un rang dont l'indice est supérieur ou égal au nombre total de fichiers reçoit une liste vide, un cas explicitement traité plutôt que laissé comme un comportement indéfini, à charge pour la boucle d'entraînement de gérer ce rang de façon appropriée (par exemple en lui faisant consommer un batch neutre plutôt que de bloquer indéfiniment la synchronisation collective des autres rangs).

Les quatre propriétés formelles établies au chapitre précédent, à savoir la couverture totale, l'absence de recouvrement, le déterminisme et l'équilibrage, découlent ici directement de deux propriétés du code : l'usage exclusif de l'opérateur modulo pour la répartition, qui garantit mécaniquement couverture et absence de recouvrement par construction de la division euclidienne, et l'ordre fixe des fichiers dans le manifeste (obtenu par un tri lexicographique systématique lors du balayage du répertoire), qui garantit que deux appels successifs avec les mêmes paramètres produisent toujours exactement la même attribution.

Cette logique d'attribution est validée par une batterie de dix cas de test, couvrant le déterminisme, la couverture totale, l'absence de recouvrement, le cas d'un rang unique, le cas où le nombre de rangs dépasse le nombre de fichiers, la distribution parfaitement équilibrée, la distribution déséquilibrée d'une unité, ainsi que trois cas de rejet explicite pour des paramètres invalides (rang négatif, rang supérieur ou égal au nombre de rangs, nombre de rangs nul). La sérialisation du manifeste lui même fait l'objet d'une batterie distincte de cinq tests, couvrant l'aller retour complet d'écriture puis de lecture, l'exactitude du total d'échantillons agrégés, la structure attendue du document JSON produit, et le rejet explicite d'un répertoire vide ou inexistant.

\bigskip
\noindent Ce chapitre a détaillé l'implémentation Python du prétraitement, depuis la configuration déclarative et la validation du schéma jusqu'à l'écriture du format binaire, en passant par le traitement d'image, la tokenisation et l'orchestration multiprocessus à mémoire bornée. Le chapitre suivant décrira le runtime C++ chargé de lire ces fichiers, d'en assembler les batches et de les transférer vers le GPU.

\begin{titlepage}
  \thispagestyle{empty} % Supprime les en-têtes/pieds pour cette page

  % LOGOS HAUT DE PAGE
  \begin{minipage}{0.48\textwidth}
  \includegraphics[width=7cm]{Images_/logo_ensam_mek.png}\\[0.2cm]
  \textbf{}
\end{minipage}
\hfill
\begin{minipage}{0.48\textwidth}
  \raggedleft
  % Pense à ajouter le logo de AI Movement dans ton dossier Images_
  \includegraphics[width=5.5cm]{Images_/aimovement_logo.png}\\[0.4cm]
  \textbf{} \\
\end{minipage}

  % ESPACE VIDE POUR PLACEMENT VERTICAL AU CENTRE
  \vspace*{5cm}

  % TITRE AU CENTRE
  \begin{center}
    \rule{\linewidth}{0.5pt}\vspace{0.7cm}

    {\Huge\bfseries\color{blue}
    Chapitre 5 :\\[0.6em] Implémentation : Runtime de Chargement C++ Haute Performance}

    \vspace{0.7cm}
    \rule{\linewidth}{0.5pt}
  \end{center}

  \vfill

\end{titlepage}

\chapter{Implémentation : Runtime de Chargement C++ Haute Performance}
\label{chap:runtime_cpp}

Le chapitre précédent a détaillé la chaîne Python qui transforme un corpus source en fichiers binaires. Ce chapitre décrit le runtime C++ qui lit ces fichiers pendant l'entraînement lui même. C'est ici que se concrétisent les décisions du cadre théorique : projection mémoire, padding dynamique, concurrence sans verrou. La vérification systématique du code source a fait apparaître plusieurs écarts entre l'intention documentée du projet et son implémentation effective. Ces écarts sont signalés explicitement au fil du texte, car ils constituent une analyse critique aussi importante que la description du fonctionnement nominal.

\section{Projection mémoire RAII}
\label{sec:mmap_raii}

La classe responsable de la projection mémoire encapsule l'ensemble du cycle de vie d'un fichier projeté dans un objet unique, non copiable mais déplaçable. Sa construction ouvre le fichier, lit sa taille par un appel système dédié, puis réalise la projection en lecture seule et en mode privé. Sur les systèmes POSIX, l'indicateur de préchargement du noyau est ajouté à la projection lorsqu'il est disponible, déclenchant la lecture anticipée de l'intégralité du fichier dès l'ouverture plutôt qu'à la demande. Un conseil d'accès séquentiel est également transmis au noyau immédiatement après la projection.

La garantie fournie par cette classe est celle de l'idiome RAII : la ressource projetée est indissociable de la durée de vie de l'objet. Le destructeur retire systématiquement la projection, y compris lorsque l'objet est détruit à la suite d'une exception levée ailleurs dans le programme, éliminant par construction toute fuite de projection mémoire. Le déplacement de l'objet est autorisé et transfère la propriété de la projection sans la dupliquer, tandis que la copie est explicitement interdite.

L'accès aux octets projetés se fait par une structure légère décrivant une plage contiguë (un pointeur et une longueur), sans posséder la mémoire qu'elle référence. Cette structure n'est valide que tant que l'objet propriétaire de la projection reste vivant, une contrainte documentée mais non vérifiée mécaniquement par le langage, et qui repose donc sur la discipline du code appelant.

Toute erreur d'ouverture, d'obtention de la taille du fichier ou de projection déclenche une exception explicite décrivant le chemin concerné, plutôt qu'un code de retour à vérifier manuellement, cohérent avec le reste de la base de code C++.

\begin{table}[htbp]
\centering
\renewcommand{\arraystretch}{1.5}
\begin{tabular}{|p{3.4cm}|p{9.6cm}|}
\hline
\textbf{Opération} & \textbf{Comportement} \\
\hline
Construction & Ouvre le fichier, lit sa taille, réalise la projection en lecture seule et mode privé, applique le préchargement noyau et le conseil d'accès séquentiel disponibles sur la plateforme \\
\hline
Destruction & Retire la projection et referme le descripteur de fichier de façon inconditionnelle \\
\hline
Accès aux données & Retourne un pointeur constant vers le début de la région projetée et sa taille totale \\
\hline
Extraction d'une plage & Retourne une plage contiguë sur un intervalle donné, avec vérification explicite de dépassement \\
\hline
\end{tabular}
\caption{Cycle de vie de la classe de projection mémoire.}
\label{tab:mmap_api}
\end{table}

Le principe théorique établi précédemment reste la justification de ce mécanisme : la projection élimine la copie que la lecture par appel système impose entre le cache de pages du noyau et un tampon applicatif. Un rappel utile ici concerne la distinction entre défaut de page mineur et majeur. Un défaut mineur survient lorsque la page demandée est déjà présente dans le cache du noyau mais pas encore associée à la table de pages du processus courant, un coût faible limité à la mise à jour de cette table. Un défaut majeur survient lorsque la page doit être effectivement chargée depuis le support de stockage, un coût correspondant à une opération d'entrée sortie complète. Le préchargement noyau appliqué à l'ouverture vise précisément à convertir par avance les défauts majeurs qui surviendraient lors des premiers accès en un unique coût amorti au moment de l'ouverture du fichier.

Une contrainte souvent oubliée mérite d'être rappelée : l'adresse de départ d'une projection doit être un multiple de la taille de page du système, une contrainte automatiquement respectée ici puisque l'intégralité du fichier est projetée depuis son octet zéro. Une implémentation qui chercherait à ne projeter qu'une plage arbitraire d'un fichier plus volumineux devrait explicitement arrondir l'offset demandé à la frontière de page inférieure la plus proche.

\section{Analyse de shards C++ et vérification d'intégrité}
\label{sec:parseur_cpp}

La vérification du code source révèle une distinction importante que la seule lecture de la spécification du format ne laisse pas deviner : il existe en réalité deux chemins de lecture distincts, tous deux capables d'interpréter le même format binaire, mais avec des garanties différentes.

Le premier chemin, le lecteur simple, ouvre le fichier par un flux d'entrée classique. Il lit l'en tête, vérifie le nombre magique de départ, lit la table des offsets, puis recalcule la somme de contrôle sur l'ensemble du contenu précédant le pied de fichier et la compare à la valeur stockée, rejetant explicitement le fichier en cas de désaccord ou de nombre magique de fin invalide. Cette vérification a lieu une seule fois, à l'ouverture. La lecture d'un échantillon individuel copie ensuite l'enregistrement correspondant dans un tampon temporaire avant de l'interpréter.

Le second chemin, interne au chargeur asynchrone décrit plus loin dans ce chapitre, ne passe pas par cette classe. Il projette directement le fichier en mémoire et relit l'en tête et la table des offsets à même la région projetée, par de simples copies de quelques octets à des positions calculées. Cette différence a une conséquence directe qu'il convient de signaler explicitement : \textbf{ce second chemin ne recalcule jamais la somme de contrôle ni ne vérifie le pied de fichier}. La vérification d'intégrité complète n'est donc garantie que par le lecteur simple, tandis que le chemin de production réellement utilisé pendant l'entraînement fait l'hypothèse implicite que les fichiers ont déjà été validés en amont, par exemple lors de la construction du manifeste au chapitre précédent, ou lors d'une inspection préalable par le lecteur simple. Ce compromis privilégie la vitesse d'ouverture (aucune relecture intégrale du fichier n'est nécessaire au démarrage) au prix d'une vérification différée qui n'est jamais rejouée par le chemin rapide lui même.

\begin{table}[htbp]
\centering
\renewcommand{\arraystretch}{1.5}
\begin{tabular}{|p{3.0cm}|p{5.0cm}|p{5.0cm}|}
\hline
\textbf{Aspect} & \textbf{Lecteur simple} & \textbf{Chemin asynchrone} \\
\hline
Accès au fichier & Flux d'entrée classique, avec recherche par position & Projection mémoire directe \\
\hline
Vérification à l'ouverture & Somme de contrôle et pied de fichier vérifiés systématiquement & Aucune vérification de somme de contrôle ni de pied de fichier \\
\hline
Lecture d'un échantillon & Copie de l'enregistrement dans un tampon temporaire avant analyse & Analyse directement depuis la région projetée, sans lecture disque supplémentaire \\
\hline
Usage prévu & Inspection, tests, accès ponctuel & Chemin de production pendant l'entraînement \\
\hline
\end{tabular}
\caption{Deux chemins de lecture distincts pour un même format binaire.}
\label{tab:deux_lecteurs}
\end{table}

Les deux chemins partagent néanmoins la même logique d'analyse d'un enregistrement d'échantillon, factorisée dans une fonction commune afin d'éviter toute divergence de comportement entre eux. Cette fonction avance un pointeur à travers le tampon fourni, avec vérification explicite de dépassement à chaque champ, lit le préfixe de dimensions lorsqu'il est présent, puis matérialise l'échantillon dans une structure qui possède ses propres tableaux. Il faut ici nuancer une affirmation répandue : la projection mémoire élimine la copie noyau vers tampon applicatif qu'imposerait une lecture classique, mais elle n'élimine pas toute copie en aval. La structure d'échantillon retournée possède ses propres tableaux, dans lesquels les octets bruts sont soit recopiés tels quels (cas des tableaux de tokens), soit transformés (cas de l'image, convertie d'entiers non signés vers des flottants normalisés). Le zéro copie annoncé au chapitre théorique désigne précisément l'absence de la double copie noyau propre à la lecture classique, non l'absence de toute écriture entre la mémoire projetée et la structure finale.

La somme de contrôle elle même n'est pas isolée dans un fichier dédié mais calculée par une table précalculée à la compilation, intégrée directement dans le lecteur simple. Le polynôme employé est identique à celui utilisé côté Python, garantissant qu'un fichier jugé intact par l'un des deux langages l'est également par l'autre. La lecture des champs entiers de l'en tête et de la table des offsets se fait systématiquement par copie mémoire explicite d'un nombre fixe d'octets vers une variable typée, plutôt que par une réinterprétation directe du pointeur, ce qui évite tout comportement indéfini lié à un alignement mémoire insuffisant sur les architectures qui l'exigent. Une version de format non reconnue déclenche un rejet explicite plutôt qu'une tentative d'interprétation optimiste.

\section{Compression et compromis débit contre taille}
\label{sec:compression_lz4}

Une option de compression peut être activée par configuration au moment de l'écriture, comme vu au chapitre précédent. Contrairement à ce que sa présence pourrait laisser supposer, seule l'interface de compression par bloc est intégrée au projet, sous la forme d'une unique implémentation source vendue directement dans l'arborescence du runtime. Il ne s'agit pas du format autodescriptif à trame complète de cette bibliothèque, mais de son mode le plus élémentaire : chaque échantillon compressé est précédé d'un entier explicite indiquant sa taille non compressée, information nécessaire à la décompression puisque le mode par bloc ne l'encode pas lui même.

À la lecture, lorsque le drapeau correspondant est positionné dans l'en tête, l'enregistrement de l'échantillon est d'abord décompressé dans un tampon temporaire dimensionné à partir de cette taille annoncée, avant d'être transmis à la fonction d'analyse commune décrite à la section précédente. Un second bit du champ de drapeaux est réservé pour une compression alternative, mais aucune implémentation de décompression correspondante n'existe actuellement dans le code : ce bit documente une extension prévue plutôt qu'une fonctionnalité disponible.

Le compromis économique de cette option se formalise simplement. Activer la compression n'apporte un gain net que si le temps de décompression ajouté par échantillon reste inférieur au temps d'entrée sortie économisé par la réduction de volume sur le support de stockage :

\begin{equation}
\label{eq:gain_compression}
t_{\text{decomp}} < t_{\text{E/S gagné}}
\end{equation}

ce qui, en introduisant le ratio de compression $r$ (volume non compressé divisé par volume compressé) et la vitesse de décompression par cœur $v_{\text{decomp}}$, se reformule comme une condition portant sur la vitesse effective du support de stockage $v_{\text{E/S}}$ :

\begin{equation}
\label{eq:condition_gain}
\frac{v_{\text{E/S}}}{v_{\text{decomp}}} < 1 - \frac{1}{r}
\end{equation}

La bibliothèque de compression retenue est spécifiquement conçue pour une décompression très rapide, de l'ordre de plusieurs gigaoctets par seconde et par cœur, au prix d'un ratio de compression modeste comparé à des algorithmes plus agressifs. Ce choix de conception traduit une priorité claire : sur un support de stockage rapide (disque à semi conducteurs local ou système de fichiers réseau à haut débit), le membre gauche de l'inégalité précédente devient rapidement défavorable à la compression, et l'option reste désactivée par défaut dans la configuration du projet. Elle ne devient pertinente que sur un support de stockage significativement plus lent que la vitesse de décompression disponible, un cas de figure que le protocole d'évaluation défini au chapitre trois permettrait de caractériser expérimentalement en variant le support de stockage cible.

\section{File sans verrou et boucle de préchargement asynchrone}
\label{sec:spsc_async}

Cette section porte sur le cœur du projet et se veut à la rigueur d'un chapitre de systèmes.

\subsection{La file sans verrou}

La structure de communication retenue est un tampon circulaire de capacité fixe, choisie librement par l'appelant et non contrainte à une puissance de deux : la progression des indices se fait par une opération modulo classique plutôt que par un masquage binaire. Une capacité utile de un élément de moins que la capacité totale est effectivement exploitable, le dernier emplacement servant à distinguer sans ambiguïté l'état plein de l'état vide.

Les deux indices de la structure, l'un lu et avancé exclusivement par le côté producteur, l'autre exclusivement par le côté consommateur, sont explicitement alignés chacun sur sa propre ligne de cache. Sans cette précaution, chaque avancée de l'un des deux indices invaliderait, au titre du protocole de cohérence de cache, la ligne entière détenue par le cœur exécutant l'autre rôle, y compris la portion de cette ligne qui ne contient pourtant aucune donnée modifiée. Ce phénomène de faux partage dégraderait sévèrement le débit de la structure malgré un coût algorithmique théoriquement nul.

L'ordre mémoire employé pour chacune de ces opérations mérite d'être justifié précisément plutôt que simplement mentionné. Le côté producteur lit sa propre position de façon détendue, car lui seul l'écrit, mais lit la position de l'autre côté avec une sémantique d'acquisition, afin d'observer sans ambiguïté tout emplacement libéré par une lecture antérieure du consommateur. Il publie ensuite le nouvel élément en écrivant sa propre position avec une sémantique de libération, garantissant que l'écriture de la donnée elle même, effectuée juste avant, est visible par le consommateur au moment où celui ci observera cette nouvelle position. Le raisonnement est strictement symétrique côté consommateur. Cette paire libération et acquisition établit une relation d'ordre garantissant qu'aucune donnée n'est lue avant d'avoir été entièrement écrite, sans qu'aucun verrou explicite ne soit nécessaire.

\subsection{Un point d'attention architectural}

La documentation de cette structure, comme son nom l'indique, suppose un unique producteur physique. Or l'examen du chargeur asynchrone révèle qu'un bassin de plusieurs threads (quatre par défaut) exécute la même boucle de traitement, et que chacun de ces threads appelle indépendamment la même opération d'insertion sur la même instance de la file, sans qu'aucun verrou supplémentaire ne protège cette opération d'insertion elle même. Deux threads du bassin peuvent ainsi, en théorie, lire simultanément la même position de dépôt avant que l'un d'eux n'ait eu le temps de la faire progresser, ce qui constituerait une situation de concurrence non définie par la structure elle même, hors du cadre pour lequel elle a été conçue et prouvée correcte. Il s'agit là d'une tension architecturale identifiée par cette relecture du code, entre l'intention documentée (un seul producteur) et l'implémentation réelle (un bassin de producteurs physiques partageant une seule opération d'insertion sans synchronisation additionnelle). Ce point constitue une piste concrète d'amélioration à traiter dans la discussion des limites de ce travail.

\subsection{Organisation de la boucle de préchargement}

Le chargeur asynchrone orchestre en réalité deux structures de communication distinctes, à ne pas confondre. La première est une file classique protégée par un verrou et une variable de condition, par laquelle un thread de répartition unique distribue des lots d'adresses d'échantillons à traiter. La seconde est la file sans verrou décrite ci-dessus, par laquelle les threads du bassin déposent les batches assemblés, consommés à leur tour par l'appelant Python.

\begin{figure}[htbp]
\centering
\begin{tikzpicture}[
  scale=0.9, transform shape,
  node distance=1.0cm,
  box/.style={draw, rectangle, rounded corners, minimum width=3.6cm, minimum height=0.9cm, align=center, font=\small},
  arr/.style={->, thick}
]
  \node[box] (order) {Ordre mélangé\\des échantillons};
  \node[box, below=of order] (dispatch) {Thread de répartition\\(file protégée par verrou)};
  \node[box, below=of dispatch, xshift=-3.2cm] (w1) {Thread\\du bassin 1};
  \node[box, below=of dispatch, xshift=3.2cm] (w2) {Thread\\du bassin N};
  \node[box, below=4.0cm of dispatch] (queue) {File sans verrou\\(profondeur de préchargement)};
  \node[box, below=of queue] (consumer) {Appelant Python\\(consommateur unique)};

  \draw[arr] (order) -- (dispatch);
  \draw[arr] (dispatch) -- (w1);
  \draw[arr] (dispatch) -- (w2);
  \draw[arr] (w1) -- (queue);
  \draw[arr] (w2) -- (queue);
  \draw[arr] (queue) -- (consumer);
\end{tikzpicture}
\caption{Deux threads ou davantage insèrent concurremment dans une même file conçue pour un seul producteur.}
\label{fig:async_arch}
\end{figure}

Chaque thread du bassin retire un lot, projette et analyse chacun des échantillons qui le composent directement depuis la mémoire projetée correspondante, assemble le batch (avec padding dynamique lorsque le format le permet), puis tente de le déposer dans la file de sortie en scrutant activement jusqu'à réussite plutôt qu'en s'endormant. Côté consommateur, l'appelant Python retire un batch en alternant tentative de retrait et attente bornée sur une variable de condition, un compromis entre réactivité et consommation processeur pendant l'attente.

Le mélange de l'ordre de présentation des échantillons est réalisé au niveau de l'ensemble des échantillons de ce rang, non au niveau d'un seul fichier à la fois, par un générateur pseudo aléatoire réamorcé de façon déterministe à partir de la somme de la graine de configuration et du numéro de l'époque courante. Lorsque la taille du tampon de mélange configuré couvre l'ensemble des échantillons, un mélange complet est réalisé. Dans le cas contraire, un algorithme de réservoir produit un mélange approché à mémoire bornée, exactement le compromis théorique annoncé au chapitre trois entre qualité du mélange et empreinte mémoire, ici concrètement instancié.

Le passage à une nouvelle époque arrête proprement les threads en cours, purge les deux files, réamorce le mélange, puis relance le bassin. Un mécanisme de point de contrôle, volontairement minimal, ne conserve que le numéro d'époque et la graine, suffisant pour restaurer intégralement l'état déterministe du chargeur sans avoir à sérialiser la position exacte de lecture dans chaque fichier.

Le débit soutenu de l'ensemble suit la relation déjà établie au chapitre trois, le minimum des taux de service des deux extrémités, et la profondeur de la file de sortie borne le temps de réponge moyen conformément à la loi de Little, une file quasiment toujours pleine ou quasiment toujours vide traduisant un déséquilibre marqué entre la vitesse de production et la vitesse de consommation.

\section{Collation dynamique de batch et normalisation différée}
\label{sec:collation_cpp}

Il convient de corriger ici une simplification commode mais inexacte : la normalisation d'entiers non signés vers des flottants n'a pas lieu au moment de la collation du batch, mais plus tôt, au moment même de l'analyse de chaque échantillon individuel, immédiatement après son extraction du tampon source. Cette normalisation est vectorisée par jeu d'instructions étendu lorsque le compilateur cible le permet, traitant huit éléments par itération au moyen d'une seule instruction combinant multiplication et addition, avec un repli scalaire pour les éléments restants sur les architectures ne supportant pas cette extension. La collation proprement dite reçoit donc des échantillons déjà normalisés et se limite à leur assemblage et à leur alignement dimensionnel.

\subsection{Deux chemins de collation}

Le premier chemin, hérité du format à dimensions fixes, suppose que tous les échantillons du batch partagent exactement les mêmes dimensions et se contente d'une copie séquentielle de chacun dans le tenseur de sortie. Le second chemin, dédié au format à dimensions variables, constitue la fonction la plus directement issue du cadre théorique de ce mémoire. Il détermine d'abord la hauteur et la largeur maximales effectivement présentes dans le batch courant, alloue un tenseur de sortie initialisé à zéro à ces dimensions maximales, puis copie chaque image ligne par ligne dans le coin supérieur gauche de son emplacement, le reste demeurant à zéro. Un masque binaire de même empreinte spatiale est simultanément construit, valant un sur les pixels réels et zéro sur les pixels de remplissage ainsi introduits, destiné à être exploité par le mécanisme d'attention du modèle pour ignorer les positions de remplissage.

\begin{table}[htbp]
\centering
\renewcommand{\arraystretch}{1.5}
\begin{tabular}{|p{4cm}|p{4.9cm}|p{4.9cm}|}
\hline
\textbf{Caractéristique} & \textbf{Chemin à dimensions fixes} & \textbf{Chemin à padding dynamique} \\
\hline
Dimensions de sortie & Identiques pour tous les batches du corpus & Recalculées pour chaque batch, au maximum du batch courant \\
\hline
Masque de remplissage & Absent & Généré, un par échantillon \\
\hline
Dimension des séquences textuelles & Longueur fixe configurée & Également fixe, non réduite dynamiquement à ce jour \\
\hline
\end{tabular}
\caption{Les deux chemins de collation exposés par la structure de batch.}
\label{tab:collation_paths}
\end{table}

Une nuance importante, absente d'une lecture rapide de la conception, doit être soulignée : seule la dimension image bénéficie effectivement de ce padding dynamique dans l'implémentation actuelle. Les quatre tableaux de tokens sont copiés à la longueur fixe configurée pour l'ensemble du corpus, pour les deux chemins de collation, sans qu'aucune réduction à la longueur réelle maximale du batch ne soit actuellement réalisée à ce niveau. Le masque d'attention associé à chaque séquence reste disponible et porte toujours l'information de longueur réelle, exploitable en aval par le modèle lui même ou par une évolution future de la collation, mais la promesse théorique d'un gain quadratique appliqué à la fois à l'image et au texte ne se vérifie, dans le code tel qu'il existe, que du côté de l'image.

Un mécanisme complémentaire regroupe les échantillons par intervalle de ratio d'aspect avant leur assemblage en batch, en un nombre configurable de tranches réparties de façon logarithmique autour d'un ratio carré. En restreignant chaque batch à des échantillons de forme comparable, ce regroupement réduit davantage l'écart entre les dimensions maximales du batch et les dimensions réelles de chaque échantillon qui le compose, complétant ainsi le padding dynamique plutôt que s'y substituant.

\section{Transfert vers le processeur graphique et mémoire épinglée}
\label{sec:gpu_transfer}

Un module distinct, compilé conditionnellement lorsqu'une boîte à outils de calcul graphique est détectée à la configuration du projet, met en œuvre l'ensemble des mécanismes classiques d'un transfert hôte vers périphérique optimisé. Une mémoire hôte épinglée (non paginable par le système d'exploitation) est allouée pour servir de zone de transit, garantissant que le contrôleur d'accès direct à la mémoire du périphérique peut transférer les données sans l'étape de copie intermédiaire vers un tampon paginable interne que le pilote insérerait sinon automatiquement. Le batch est transféré champ par champ (image puis chacun des quatre tableaux de tokens) vers cette zone de transit, chaque copie étant suivie d'un transfert asynchrone vers une mémoire allouée sur le périphérique. Les pointeurs de périphérique obtenus sont ensuite exposés selon le protocole d'échange de tenseurs entre bibliothèques, sous la forme d'une structure gérée dotée d'un destructeur personnalisé qui libère la mémoire de périphérique lorsque la propriété lui en a été explicitement transférée, permettant à un cadriciel d'apprentissage de consommer ce pointeur sans copie supplémentaire tout en garantissant sa libération correcte.

Un examen attentif révèle toutefois une réutilisation d'une seule et même zone de transit pour les cinq champs successifs du batch, ce qui impose une synchronisation explicite du flux après chaque champ avant de réutiliser cette zone pour le champ suivant. Les cinq transferts sont donc, en pratique, sérialisés les uns par rapport aux autres malgré l'usage d'une primitive de copie asynchrone : le gain obtenu porte sur l'élimination de la copie intermédiaire du pilote pour chaque transfert individuel, non sur un recouvrement mutuel des cinq transferts entre eux.

Un second examen, plus important encore, doit être signalé sans détour : \textbf{ce module n'est actuellement invoqué par aucun point d'entrée exposé à Python}. La méthode réellement accessible depuis l'interface Python, qui combine l'obtention du prochain batch et sa conversion vers le cadriciel d'apprentissage, procède différemment : chaque tableau du batch est d'abord copié vers un tableau numérique standard, puis, si le cadriciel d'apprentissage est disponible dans l'environnement, converti en tenseur par la fonction d'importation directe de ce cadriciel, et déplacé vers le périphérique cible par l'appel de conversion générique de ce même cadriciel lorsqu'un périphérique autre que le processeur central est demandé. Cette dernière étape correspond à un transfert hôte vers périphérique ordinaire, non épinglé, non asynchrone au sens du module décrit ci-dessus.

\begin{table}[htbp]
\centering
\renewcommand{\arraystretch}{1.5}
\begin{tabular}{|p{4.0cm}|p{5.0cm}|p{4.0cm}|}
\hline
\textbf{Étape} & \textbf{Module de transfert optimisé} & \textbf{Chemin réellement exposé} \\
\hline
Zone de transit & Mémoire hôte épinglée, réutilisée par champ & Tableau numérique standard, une copie par champ \\
\hline
Transfert vers le périphérique & Copie asynchrone suivie d'une synchronisation par champ & Appel de conversion générique du cadriciel d'apprentissage \\
\hline
Remise au cadriciel & Protocole d'échange de tenseurs, sans copie supplémentaire & Conversion depuis le tableau numérique, une copie supplémentaire \\
\hline
Statut & Implémenté, non relié à l'interface exposée & Chemin effectivement emprunté à chaque appel \\
\hline
\end{tabular}
\caption{Le module de transfert optimisé existe mais n'est pas encore branché sur l'interface exposée.}
\label{tab:gpu_gap}
\end{table}

Cet écart entre une infrastructure de transfert avancée et une interface exposée qui ne l'exploite pas encore constitue une perspective d'amélioration directe et clairement identifiée pour la suite de ce travail, plutôt qu'un défaut de conception : les deux briques existent séparément et leur intégration ne demande pas de nouvelle conception, seulement le branchement de l'une sur l'autre au niveau de la couche d'exposition.

Le coût théorique de chacune des deux stratégies reste celui établi précédemment. Un transfert non épinglé impose une copie hôte intermédiaire avant le transfert par accès direct à la mémoire, soit deux copies au total, tandis qu'un transfert depuis une zone épinglée n'impose qu'une seule copie effective, celle du transfert par accès direct à la mémoire lui même, la première zone étant déjà dans un état compatible avec ce transfert.

\section{Coordination distribuée et interface exposée à Python}
\label{sec:coordination_bindings}

\subsection{Deux niveaux de partitionnement distincts}

Il est tentant, à la lecture du plan de ce chapitre, de supposer que le partitionnement distribué réalisé côté C++ reproduit exactement celui réalisé côté Python au chapitre précédent. L'examen du code montre qu'il s'agit en réalité de deux mécanismes distincts, opérant à des granularités différentes et sans lien direct entre eux dans le code lui même.

Le mécanisme Python, déjà décrit, répartit des fichiers entiers entre rangs par une répartition cyclique fondée sur l'indice du fichier dans le manifeste. Le mécanisme C++ décrit ici répartit, à l'intérieur d'un ensemble de fichiers déjà fourni au chargeur (quelle que soit la façon dont cet ensemble a été déterminé en amont), les indices d'échantillons individuels selon deux stratégies au choix : une répartition en blocs contigus, privilégiant la localité de lecture au prix d'un déséquilibre potentiel si les tailles d'échantillons varient fortement, ou une répartition entrelacée, répartissant les indices en alternance et privilégiant l'équilibrage au prix d'une localité de lecture moindre. Rien dans le code ne relie ce second mécanisme au manifeste JSON produit côté Python : c'est à l'appelant de la bibliothèque de décider, en amont, quels fichiers transmettre au chargeur, le manifeste et l'attribution par rang restant un outil Python que le runtime C++ n'interroge jamais directement.

Les deux mécanismes partagent néanmoins les mêmes exigences formelles de couverture totale et d'absence de recouvrement, chacun disposant de sa propre fonction de vérification testant cette propriété par construction plutôt que par un partage de code entre les deux langages.

\subsection{Interface exposée à Python}

L'interface Python de ce runtime expose la structure de configuration, le chargeur asynchrone lui même selon le protocole d'itération natif, ainsi que les utilitaires de partitionnement, de vérification et de regroupement par ratio d'aspect. Chaque appel bloquant traversant cette interface (l'obtention du prochain batch, sous ses différentes formes d'accès) retient, par défaut, l'interpréteur pendant toute sa durée d'exécution, aucune déclaration explicite de relâchement n'étant présente dans le code d'exposition. Ceci mérite d'être signalé avec précision : l'attente potentiellement bloquante à l'intérieur de l'obtention du prochain batch s'exécute donc verrou tenu, alors que les threads internes du bassin de préchargement, qui n'appellent eux mêmes jamais l'interface Python, continuent de progresser indépendamment de ce verrou, celui ci ne les concernant pas. La conséquence pratique est que le thread Python appelant reste indisponible pour toute autre tâche pendant l'attente, une limitation mineure dans un scénario d'entraînement à thread Python unique, mais qui mériterait d'être levée par un relâchement explicite si l'interface devait un jour être appelée depuis plusieurs threads Python concurrents.

La chaîne de compilation du module d'extension repose entièrement sur un outil de configuration multiplateforme plutôt que sur un script d'empaquetage Python séparé, ce même outil détectant automatiquement la disponibilité d'une boîte à outils de calcul graphique et l'extension vectorielle du processeur hôte, activant conditionnellement les chemins de code correspondants sans intervention manuelle.

\bigskip
\noindent Ce chapitre a détaillé le runtime C++ et, par une lecture attentive du code plutôt qu'une simple paraphrase de son intention documentée, a mis en évidence plusieurs écarts significatifs entre la conception visée et l'implémentation actuelle : l'absence de vérification d'intégrité sur le chemin de production, une file conçue pour un seul producteur mais alimentée par plusieurs, un padding dynamique qui ne couvre à ce jour que la dimension image, et un module de transfert optimisé développé mais non encore relié à l'interface exposée. Ces constats nourriront directement la discussion des limites et des perspectives de ce travail, développée dans les chapitres suivants.


\end{document}